\documentclass[../../main.tex]{subfiles}% 注意这里的文件路径不能用 ./main.tex ,否则用latexmk编译子文件会报错
\graphicspath{{\subfix{./image/}}{\subfix{../../image/}}} % 指定图片引用路径,后续可以直接使用图片文件名
% 如果图片存放在上级目录,则使用 ../../image/ 作为备用路径

% 例如:
% \begin{figure}[H]
% \centering
% \includegraphics[width=0.3\textwidth]{图.png}
% \caption{}
% \label{figure:图}
% \end{figure}
% 注意:上述\label{}一定要放在\caption{}之后,否则引用图片序号会只会显示??.

\begin{document}

\section{共轭空间、弱收敛、自反空间}

\subsection{共轭空间的表示及应用}

\begin{theorem}\label{theorem:共轭空间的完备性}
设$\mathscr{X}$是一个$B^*$空间，$\mathscr{X}$上的所有连续线性泛函全体$\mathscr{X}^*$，按范数
\begin{align*}
\|f\|=\sup\limits_{\|x\|=1}|f(x)|
\end{align*}
构成一个$B$空间，称为$\mathscr{X}$的\textbf{共轭空间}或\textbf{对偶空间}。
\end{theorem}
\begin{proof}
$\mathscr{X}^*$的完备性可由$\mathscr{X}^*=\mathscr{L}(\mathscr{X},\mathbb{K})$和\refthe{theorem:有界线性算子全体构成的空间完备}直接导出.
\end{proof}

\begin{example}
定义空间
\begin{align*}
BV[0,1]\triangleq\left\{g:[0,1]\to\mathbb{C}\,\middle|\,
\begin{aligned}
&g(0)=0,\\
&g(t)=g(t+0)\ (\forall t\in(0,1)),\\
&\mathrm{var}(g)<\infty
\end{aligned}
\right\},
\end{align*}
其中变差$\mathrm{var}(g)=\sup\sum\limits_{j=0}^{n-1}|g(t_{j+1})-g(t_j)|$，上确界取遍$[0,1]$的所有分割$\Delta:0=t_0<t_1<\cdots<t_n=1$。在$BV[0,1]$上赋范$\|g\|_v=\mathrm{var}(g)\ (\forall g\in BV[0,1])$，则$BV[0,1]$为$B$空间，且
\begin{align}
C[0,1]^*=BV[0,1].\label{eq:cconj}
\end{align}
\end{example}
\begin{remark}
有界变差函数可表示为两个单调递增函数的差$g=g_1-g_2$，调整函数取值可使其右连续，且不改变Stieltjes积分。映射$g\mapsto\int_{0}^{1}\varphi(t)\mathrm{d}g(t)$是$BV[0,1]\to C[0,1]^*$的等距同构。
\end{remark}
\begin{proof}
对任意$\varphi\in C[0,1],g\in BV[0,1]$，Stieltjes积分定义为
\begin{align*}
\int_{0}^{1}\varphi(t)\mathrm{d}g(t)=\lim\limits_{\|\Delta\|\to0}\sum\limits_{j=0}^{n-1}\varphi(t_j^*)[g(t_{j+1})-g(t_j)],
\end{align*}
其中$\|\Delta\|=\max\limits_{1\leqslant j\leqslant n}|t_j-t_{j-1}|$，$t_j^*\in[t_j,t_{j+1}]$。

对任意$g\in BV[0,1]$，其诱导$C[0,1]$上的连续线性泛函
\begin{align*}
\varphi\mapsto\langle f,\varphi\rangle=\int_{0}^{1}\varphi(t)\mathrm{d}g(t),
\end{align*}
满足
\begin{align}
\|f\|\leqslant\int_{0}^{1}|\mathrm{d}g(t)|=\mathrm{var}(g)=\|g\|_v.\label{eq:bstieltnorm}
\end{align}

反之，对任意$f\in C[0,1]^*$，由\hyperref[theorem:泛函分析--(哈恩-巴拿赫)Hahn-Banach定理]{Hahn-Banach定理}，将$f$保范延拓为$\widetilde{f}\in L^\infty[0,1]^*$，满足$\langle\widetilde{f},\chi_{\{0\}}\rangle=0$，$\langle\widetilde{f},\varphi\rangle=\langle f,\varphi\rangle\ (\forall\varphi\in C[0,1])$且$\|\widetilde{f}\|=\|f\|$。

令
\begin{align*}
\begin{cases}
g(s)=\langle\widetilde{f},\chi_{(0,s]}\rangle,&0<s\leqslant1,\\
g(0)=0,
\end{cases}
\end{align*}
下证$g\in BV[0,1]$，且满足对应关系。

任取$[0,1]$的分割$\Delta:0=t_0<t_1<\cdots<t_n=1$，记$\omega_k\triangleq g(t_{k+1})-g(t_k)$，
\[
\lambda_k\triangleq
\begin{cases}
\overline{\omega_k}/|\omega_k|,&\omega_k\neq0,\\
0,&\text{其他},
\end{cases}
\quad
h_\Delta(t)\triangleq\sum\limits_{k=0}^{n-1}\lambda_k\chi_{(t_k,t_{k+1}]}(t).
\]
则
\begin{align*}
\sum\limits_{k=0}^{n-1}|\omega_k|=\sum\limits_{k=0}^{n-1}\lambda_k\omega_k=\langle\widetilde{f},h_\Delta\rangle
\leqslant\|\widetilde{f}\|\cdot\|h_\Delta\|_{L^\infty[0,1]}\leqslant\|\widetilde{f}\|=\|f\|,
\end{align*}
故$g\in BV[0,1]$，且$\|g\|_v\leqslant\|f\|$。

再证$\langle f,\varphi\rangle=\int_{0}^{1}\varphi(t)\mathrm{d}g(t)\ (\forall\varphi\in C[0,1])$。对任意$\varepsilon>0$，取分割$\Delta$使得
\begin{align*}
|\varphi(t)-\varphi(t')|<\frac{\varepsilon}{2\|\widetilde{f}\|},\quad\forall t,t'\in[t_j,t_{j+1}],\ j=0,1,\cdots,n-1,
\end{align*}
且
\begin{align*}
\left|\int_{0}^{1}\varphi(t)\mathrm{d}g(t)-\sum\limits_{j=0}^{n-1}\varphi(t_j)(g(t_{j+1})-g(t_j))\right|<\frac{\varepsilon}{2}.
\end{align*}
令$\varphi_\Delta\triangleq\sum\limits_{j=0}^{n-1}\varphi(t_j)\chi_{(t_j,t_{j+1}]}+\varphi(0)\chi_{\{0\}}$，则
\begin{align*}
\left|\langle f,\varphi\rangle-\int_{0}^{1}\varphi(t)\mathrm{d}g(t)\right|
&\leqslant\left|\langle f,\varphi\rangle-\langle\widetilde{f},\varphi_\Delta\rangle\right|+\left|\langle\widetilde{f},\varphi_\Delta\rangle-\int_{0}^{1}\varphi(t)\mathrm{d}g(t)\right|\\
&\leqslant\|\widetilde{f}\|\cdot\|\varphi-\varphi_\Delta\|_{L^\infty[0,1]}+\left|\sum\limits_{j=0}^{n-1}\varphi(t_j)(g(t_{j+1})-g(t_j))-\int_{0}^{1}\varphi(t)\mathrm{d}g(t)\right|\\
&<\frac{\varepsilon}{2}+\frac{\varepsilon}{2}=\varepsilon.
\end{align*}
由$\varepsilon$的任意性，等式成立。
\end{proof}

\begin{theorem}[Riesz表示定理(连续函数空间)]\label{theorem:泛函分析--Riesz表示定理连续函数空间}
若$M$是一个Hausdorff紧空间，则$\forall f\in C(M)^*$，存在唯一的复值Baire测度（完全可加的集函数）$\mu$，适合$|\mu|(M)<\infty$，满足
\begin{align}
\langle f,\varphi\rangle = \int_{M}\varphi(m)\mathrm{d}\mu \quad (\forall \varphi\in C(M)).
\end{align}
\end{theorem}

\begin{lemma}\label{lemma:泛函分析--引理22421}
对$\forall \mu\in\mathscr{M}(K)$，设
\begin{align*}
\wideparen{\mu}(w)=\int_{K}\frac{\mathrm{d}\mu(z)}{w-z},
\end{align*}
那么对$\forall R>0$，有$\wideparen{\mu}\in L^1(B_R)$（其中$B_R$是中心在原点的半径为$R$的圆），且$\wideparen{\mu}$在$\mathbb{C}_\infty\setminus K$上解析，还满足$\wideparen{\mu}(\infty)=0$.
\end{lemma}
\begin{proof}
\begin{enumerate}[(1)]
\item 证$\wideparen{\mu}\in L^1(B_R)$. 由定义
\[
|\wideparen{\mu}(w)|\leqslant\int_{K}\frac{\mathrm{d}|\mu|(z)}{|w-z|},
\]
从而
\begin{align*}
\int_{B_R}|\wideparen{\mu}(w)|\mathrm{d}x\mathrm{d}y&\leqslant\int_{B_R}\int_{K}\frac{\mathrm{d}|\mu|(z)}{|w-z|}\mathrm{d}x\mathrm{d}y
=\int_{K}\int_{B_R}\frac{\mathrm{d}x\mathrm{d}y}{|w-z|}\mathrm{d}|\mu|(z)\\
&\leqslant\int_{K}\int_{B(z,\rho)}\frac{\mathrm{d}x\mathrm{d}y}{|w-z|}\mathrm{d}|\mu|(z)\leqslant2\pi\rho|\mu|(K)<\infty,
\end{align*}
其中$\rho>R+\max\limits_{z\in K}|z|$.

\item 证$\wideparen{\mu}$在$\mathbb{C}\setminus K$解析. 因$K$紧，可在积分号下求微商.

\item 证$\wideparen{\mu}$在$\{\infty\}$解析，且$\wideparen{\mu}(\infty)=0$. 因$K$紧，当$|w|\to\infty$时，在积分号下取极限得$\wideparen{\mu}(w)\to0$，故$\infty$是可去奇点.
\end{enumerate}
\end{proof}

\begin{theorem}[Runge定理]
设$K$是复平面$\mathbb{C}$上的一个紧子集，记$\mathbb{C}_\infty=\mathbb{C}\cup\{\infty\}$。又设$E$是$\mathbb{C}_\infty\setminus K$中的一个子集，它与$\mathbb{C}_\infty\setminus K$的每一个连通分量都相交。若$f$是$K$的一个邻域内的任意解析函数，则必有有理函数列$f_n$，其极点都在$E$内，使得$f_n$在$K$上一致收敛到$f$。
\end{theorem}

\begin{proof}
引用$K$上的连续函数空间$C(K)$，记$R(K,E)$为极点在$E$内的有理函数集在$C(K)$中的闭包，则$R(K,E)$是$C(K)$的闭线性子空间。为证$f\in R(K,E)$，只需证：对$\forall F\in C(K)^*$，
\[
F(g)=0\quad (\forall g\in R(K,E))\implies F(f)=0.
\]
由\hyperref[theorem:泛函分析--Riesz表示定理连续函数空间]{Riesz表示定理(连续函数空间)}，$C(K)^*$由$K$上的完全可加的复值测度$\mathscr{M}(K)$组成，故只需证：对$\forall \mu\in\mathscr{M}(K)$，
\begin{align}
\int_{K}g\mathrm{d}\mu=0\ (\forall g\in R(K,E))\implies \int_{K}f\mathrm{d}\mu=0. \label{eq:::HIOUR32s2.5.15}
\end{align}
由\reflem{lemma:泛函分析--引理22421}可得
\begin{align}
\left(\frac{\mathrm{d}}{\mathrm{d}w}\right)^n\wideparen{\mu}(w_0)=n!\int_{K}(z-w_0)^{-n+1}\mathrm{d}\mu(z)\quad (\forall w_0\in\mathbb{C}\setminus K), \label{eq:::HIOUR32s2.5.16}
\end{align}
而在$\infty$点附近它有展开
\begin{align}
\wideparen{\mu}(w)=-\frac{1}{w}\sum\limits_{n=0}^{\infty}\int_{K}\left(\frac{z}{w}\right)^n\mathrm{d}\mu(z)=-\sum\limits_{n=0}^{\infty}\frac{a_n}{w^{n+1}}, \label{eq:::HIOUR32s2.5.17}
\end{align}
其中$a_n=\int_{K}z^n\mathrm{d}\mu(z)$.

若$\mu\in\mathscr{M}(K)$使得
\begin{align}
\int_{K}g(z)\mathrm{d}\mu(z)=0, \label{eq:::HIOUR32s2.5.18}
\end{align}
其中$g$是极点在$E$内的有理函数，则
\begin{align}
\wideparen{\mu}(w)=0\quad (\forall w\in\mathbb{C}_\infty\setminus K). \label{eq:::HIOUR32s2.5.19}
\end{align}

事实上，$\forall w_0\in E$，设$\Omega(w_0)$是$\mathbb{C}_\infty\setminus K$中含$w_0$的分量，由$E$的假设，
\begin{align}
\mathbb{C}_\infty\setminus K=\bigcup\limits_{w_0\in E}\Omega(w_0). \label{eq:::HIOUR32s2.5.20}
\end{align}
若$w_0\neq\infty$，由\eqref{eq:::HIOUR32s2.5.16}式与假设\eqref{eq:::HIOUR32s2.5.18}式，$\wideparen{\mu}$在$w_0$的各阶导数均为$0$，故$\wideparen{\mu}$在$\Omega(w_0)$内恒为$0$；若$w_0=\infty$，则由\eqref{eq:::HIOUR32s2.5.17}式与假设\eqref{eq:::HIOUR32s2.5.18}式，$\wideparen{\mu}$在$\Omega(w_0)$内也恒为$0$. 于是由\eqref{eq:::HIOUR32s2.5.20}式即得\eqref{eq:::HIOUR32s2.5.19}式.

考察在$K$的某邻域$G$内解析的任意函数$f$，存在含于$G\setminus K$内的折线$\gamma_1,\gamma_2,\cdots,\gamma_n$使得
\begin{align*}
f(z)=\sum\limits_{k=1}^{n}\frac{1}{2\pi\mathrm{i}}\int_{\gamma_k}\frac{f(w)}{w-z}\mathrm{d}w.
\end{align*}
由Fubini定理，
\begin{align}
\int_{K}f(z)\mathrm{d}\mu(z)=\sum\limits_{k=1}^{n}\frac{1}{2\pi\mathrm{i}}\int_{K}\int_{\gamma_k}\frac{f(w)}{w-z}\mathrm{d}w\mathrm{d}\mu(z)
=\sum\limits_{k=1}^{n}\frac{1}{2\pi\mathrm{i}}\int_{\gamma_k}f(w)\wideparen{\mu}(w)\mathrm{d}w=0, \label{eq:::HIOUR32s2.5.21}
\end{align}
因$\gamma_k\subseteq\mathbb{C}\setminus K$，故$\wideparen{\mu}(w)=0$于$\gamma_k$上. 于是由\eqref{eq:::HIOUR32s2.5.18}式推出\eqref{eq:::HIOUR32s2.5.21}式，即证得\eqref{eq:::HIOUR32s2.5.15}式.
\end{proof}

\begin{definition}
设 $\mathscr{X}$ 是 $B^*$ 空间，其共轭空间 $\mathscr{X}^*$ 是 $B$ 空间，$\mathscr{X}^{**} = (\mathscr{X}^*)^*$ 称为 $\mathscr{X}$ 的\textbf{第二共轭空间}。
\end{definition}

\begin{theorem}\label{theorem:原空间可等距嵌入第二共轭空间}
设 $\mathscr{X}$ 是 $B^*$ 空间，对$\forall x \in \mathscr{X}$,定义 $X : \mathscr{X}^* \to \mathbb{K}$如下：
\[
X(f) = \langle f, x \rangle\triangleq f(x) \quad (\forall f \in \mathscr{X}^*),
\]
则$X \in \mathscr{X}^{**}$.再定义映射
\begin{align*}
T : \mathscr{X} \to \mathscr{X}^{**},\,\,x\mapsto X.
\end{align*}
称映射$T$为 $\mathscr{X}$到$\mathscr{X}^{**}$ 的\textbf{自然映射}或\textbf{自然嵌入}.则$T$是$\mathscr{X}$到$\mathscr{X}^{**}$的等距嵌入,即$T$是$\mathscr{X}$到$R(T)$的等距同构。并且
\begin{align}\label{eq::jjjjjjjjjjjiiiiiiiiiiijs}
Tx\left( f \right) =\left<Tx,f \right> =\left< f,x \right> =f\left( x \right) \quad \left( \forall f\in \mathscr{X} ^*,\forall x\in \mathscr{X} \right) .
\end{align}
特别地，若$\mathscr{X}$到$\mathscr{X}^{**}$的自然映射$T$是满射的，则称$\mathscr{X}$是\textbf{自反的}，记作$\mathscr{X}=\mathscr{X}^{**}$.也即对$\forall g\in X^{**}$,都存在$x\in X$,使得$g=Tx$,进而
\begin{align*}
g(f)=Tx(f)=f(x),\quad \forall f\in X^*.
\end{align*}
\end{theorem}
\begin{remark}
有时对$x$与$X$不加区别，简记$\mathscr{X}\subseteq\mathscr{X}^{**}$.也可以用\eqref{eq::jjjjjjjjjjjiiiiiiiiiiijs}式直接定义$T$.
\end{remark}
\begin{proof}
由定义易知$X$ 是 $\mathscr{X}^*$ 上的线性泛函，且满足
\begin{align}\label{eq::IOUHJIOW22}
|X(f)| \leqslant \|f\| \cdot \|x\| \quad (\forall f \in \mathscr{X}^*),
\end{align}
故 $X$ 连续，即 $X \in \mathscr{X}^{**}$。且
\begin{align}\label{eq::--hruhiow3jio2}
\left\| X \right\| \leqslant \left\| x \right\| .
\end{align}
由\eqref{eq::IOUHJIOW22}式知映射$T:x\mapsto X$是$\mathscr{X}$到$\mathscr{X}^{**}$的连续嵌入。$\forall \alpha,\beta\in\mathbb{C},x,y\in\mathscr{X}$，记$X=Tx,Y=Ty$，则
\begin{align*}
T(\alpha x+\beta y)(f)&=f(\alpha x+\beta y)
=\alpha f(x)+\beta f(y)=\alpha X(f)+\beta Y(f)\\
&=(\alpha X+\beta Y)(f)=(\alpha Tx+\beta Ty)(f)\quad (\forall f\in\mathscr{X}^*).
\end{align*}
故$T$是线性同构。
再由\refcor{corollary:泛函分析--推论2.4.6}，$\exists f\in\mathscr{X}^*$，使得$\|f\|=1$，且$\langle f,x\rangle=\|x\|$，于是
\begin{align}
\|x\|=X(f)\leqslant\|X\|\cdot\|f\|=\|X\|. \label{eq:::HIOUR32s2.5.24}
\end{align}
结合\eqref{eq::--hruhiow3jio2}式与\eqref{eq:::HIOUR32s2.5.24}式知$\|X\|=\|x\|$，故$T$是等距的.
\end{proof}

\begin{proposition}
设 \(H\) 是一个Hilbert空间，则 \(H\) 是自反的，并且存在从 \(H\) 到\(H^{**}\) 的一个等距同构,即$H\cong H^{**}.$
\end{proposition}
\begin{proof}

由\cref{theorem:原空间可等距嵌入第二共轭空间},定义自然嵌入映射 \(J: H \to H^{**}\) 如下：对于任意 \(x \in H\)，令 \(J(x) \in H^{**}\) 满足
\begin{align}
J(x)(f) = f(x), \quad \forall f \in H^{*}.
\end{align}
显然，\(J\) 是一个线性保距映射，即 \(\|J(x)\|_{H^{**}} = \|x\|_{H}\)。

由\hyperref[theorem:Riesz表示定理-泛函]{Riesz表示定理}知，存在保距同构 \(\Phi: H \to H^{*}\)，使得 \(\Phi(x)(y) = (y, x)\)。故$H\cong H^*$,于是$H^*$也是Hilbert空间.同理由\hyperref[theorem:Riesz表示定理-泛函]{Riesz表示定理}知，存在保距同构 \(\Psi: H^{*} \to H^{**}\).因此复合映射 \(\Psi \circ \Phi: H \to H^{**}\) 也是一个保距同构.
现证明自然映射 \(J\) 是满射。任取 \(F \in H^{**}\)，由于 \(H^{**}\) 是 \(H^*\) 的对偶空间，由 \hyperref[theorem:Riesz表示定理(Hilbert空间)]{Riesz表示定理}在 \(H^*\) 上的应用，存在唯一的 \(f \in H^{*}\)，使得
\begin{align}
F(g) = \langle g, f \rangle_{H^*}, \quad \forall g \in H^{*}.
\end{align}
同时，由\hyperref[theorem:Riesz表示定理(Hilbert空间)]{Riesz表示定理}对 \(H\) 的应用，存在唯一的 \(x \in H\)，使得
\begin{align*}
f(y) = (y, x)_H, \quad \forall y \in H.
\end{align*}
于是对于任意 \(g \in H^{*}\)，由 \(H\) 上的\hyperref[theorem:Riesz表示定理(Hilbert空间)]{Riesz表示定理}，存在 \(y_g \in H\) 使得 \(g(y) = (y, y_g)\)，于是有
\begin{align*}
F(g) = \langle g, f \rangle_{H^*} = (y_g, x)_H = g(x) = J(x)(g).
\end{align*}
这说明 \(F = J(x)\)，即 \(J\) 是满射，从而 \(J\) 是一个等距同构。因此 \(H\) 是自反的，即 \(H \cong H^{**}\).
\end{proof}

\begin{proposition}\label{proposition:Lp对偶与自反性}
设 $(\Omega,\mathscr{B},\mu)$ 是 $\sigma$-有限的完全可加测度空间，$1\leqslant p<\infty$，$q$ 是 $p$ 的共轭数，即当 $p>1$ 时 $\frac{1}{p}+\frac{1}{q}=1$，当 $p=1$ 时 $q=\infty$。则在等距同构意义下，有
\begin{enumerate}[(1)]
\item $(L^p(\Omega,\mathscr{B},\mu))^*=L^q(\Omega,\mathscr{B},\mu)$，且当 $1<p<\infty$ 时 $L^p(\Omega,\mathscr{B},\mu)$ 是自反的；

\item $(L^1(\Omega,\mathscr{B},\mu))^*=L^\infty(\Omega,\mathscr{B},\mu)$，但 $L^1(\Omega,\mathscr{B},\mu)$ 不自反；

\item $L^\infty(\Omega,\mathscr{B},\mu)$ 不自反，且 $(L^\infty(\Omega,\mathscr{B},\mu))^*\neq L^1(\Omega,\mathscr{B},\mu)$，即 $(L^\infty(\Omega,\mathscr{B},\mu))^*$ 严格包含 $L^1(\Omega,\mathscr{B},\mu)$。
\end{enumerate}
\end{proposition}
\begin{proof}

先证 $(1)$。对任意 $g\in L^q(\Omega,\mathscr{B},\mu)$，定义泛函
\begin{align*}
T_g(f)\triangleq\int_{\Omega} f(x)g(x)\mathrm{d}\mu,\quad f\in L^p(\Omega,\mathscr{B},\mu).
\end{align*}
由\hyperref[Real Analysis-theorem:实变函数--Hölder不等式]{Hölder不等式}，$T_g\in (L^p(\Omega,\mathscr{B},\mu))^*$，且 $\|T_g\|\leqslant\|g\|_{L^q}$。反之，任取 $F\in (L^p(\Omega,\mathscr{B},\mu))^*$，由 \hyperref[theorem:泛函分析--Radon-Nikodym定理]{Radon-Nikodym定理}，存在 $g\in L^q(\Omega,\mathscr{B},\mu)$，使得
\begin{align*}
F(f)=\int_{\Omega} f(x)g(x)\mathrm{d}\mu,\quad \forall f\in L^p(\Omega,\mathscr{B},\mu),
\end{align*}
且 $\|g\|_{L^q}=\|F\|$。故 $(L^p(\Omega,\mathscr{B},\mu))^*=L^q(\Omega,\mathscr{B},\mu)$。

当 $1<p<\infty$ 时，由于 $(L^p)^*=L^q$ 且 $(L^q)^*=L^p$，故 $(L^p)^{**}=L^p$，并且典则映射 $J:L^p\to (L^p)^{**}$ 是满射，所以 $L^p$ 自反。

再证 $(2)$。当 $p=1$ 时，同理可得 $(L^1)^*=L^\infty$。为证 $L^1$ 不自反，只需证明 $(L^\infty)^*\neq L^1$。

下面证明 $(3)$ 中 $(L^\infty)^*\neq L^1$。仍以 $\Omega=[0,1]$，$\mu$ 为 Lebesgue 测度为例，一般 $\sigma$-有限非原子测度空间情形完全类似。

对任意 $g\in L^1[0,1]$，定义泛函
\begin{align*}
T_g(h)=\int_0^1 h(x)g(x)\mathrm{d}x,\quad h\in L^\infty[0,1].
\end{align*}
则 $T_g\in (L^\infty[0,1])^*$，且 $\|T_g\|=\|g\|_{L^1}$，故 $L^1[0,1]$ 可等距嵌入 $(L^\infty[0,1])^*$。

考虑闭子空间 $C[0,1]\subseteq L^\infty[0,1]$，在其上定义有界线性泛函
\begin{align*}
\delta_0(f)=f(0),\quad f\in C[0,1].
\end{align*}
由\hyperref[theorem:泛函分析--(哈恩-巴拿赫)Hahn-Banach定理]{Hahn-Banach定理}，存在 $T\in (L^\infty[0,1])^*$，使得 $T(f)=f(0)$ 对一切 $f\in C[0,1]$ 成立。

假设存在 $g\in L^1[0,1]$，使得
\begin{align*}
T(h)=\int_0^1 h(x)g(x)\mathrm{d}x,\quad \forall h\in L^\infty[0,1].
\end{align*}
对每个 $m\in\mathbb{N}$，令
\begin{align*}
f_m(x)=\max\{0,1-mx\},\quad x\in[0,1].
\end{align*}
则 $f_m\in C[0,1]$，且
\begin{align*}
f_m(0)=1,\quad 0\leqslant f_m(x)\leqslant 1,\quad \operatorname{supp} f_m\subseteq\left[0,\frac{1}{m}\right].
\end{align*}
于是 $f_m\to 0$ a.e. 于 $[0,1]$。又因为 $|f_m(x)g(x)|\leqslant |g(x)|$ 且 $g\in L^1[0,1]$，由 \hyperref[Real Analysis-theorem:控制收敛定理]{Lebesgue控制收敛定理}得
\begin{align*}
\lim\limits_{m\to\infty}\int_0^1 f_m(x)g(x)\mathrm{d}x=0.
\end{align*}
但另一方面 $T(f_m)=f_m(0)=1$，矛盾！因此 $T$ 不能由任何 $L^1[0,1]$ 函数表示，故 $(L^\infty[0,1])^*\neq L^1[0,1]$。

因此 $L^1$ 不自反；若 $L^\infty$ 自反，则其对偶 $L^1$ 也自反，矛盾！故 $L^\infty$ 也不自反。
\end{proof}


\subsection{共轭算子}

\begin{definition}[共轭算子]
设$\mathscr{X},\mathscr{Y}$是$B^*$空间，算子$T\in\mathscr{L}(\mathscr{X},\mathscr{Y})$。算子$T^*:\mathscr{Y}^*\to\mathscr{X}^*$称为是$T$的\textbf{共轭算子}，若满足
\begin{align}\label{eq:JHIODWI231122888}
f(Tx)=(T^*f)(x)\quad (\forall f\in\mathscr{Y}^*,\forall x\in\mathscr{X}).
\end{align}
也即
\begin{align*}
\left< f,Tx \right> =\left< T^*f,x \right> \quad (\forall f\in \mathscr{Y} ^*,\forall x\in \mathscr{X} ).
\end{align*}
\end{definition}

\begin{proposition}\label{proposition:Banach空间共轭算子必存在唯一}
设$\mathscr{X},\mathscr{Y}$是$B^*$空间，则$\forall T\in\mathscr{L}(\mathscr{X},\mathscr{Y})$，$T^*$唯一存在，且
\begin{align}\label{eq:::HIOUR32s2.5.26}
T^*\in\mathscr{L}(\mathscr{Y}^*,\mathscr{X}^*),\quad \|T^*\|\leqslant\|T\|.
\end{align}
\end{proposition}
\begin{proof}
事实上，对$\forall f\in\mathscr{Y}^*$，令$g(x)=f(Tx)\ (\forall x\in\mathscr{X})$，则$g$线性,且由\refpro{proposition:有界线性算子的基本不等式}知
\[
|g(x)|\leqslant\|f\|\cdot\|T\|\cdot\|x\|\quad (\forall x\in\mathscr{X}).
\]
故$g\in\mathscr{X}^*$.令$T^*:\mathscr{Y}^*\to\mathscr{X}^*,f\mapsto g$线性,则由定义知$T^*$是$T$的共轭算子,且
\begin{align}
\|T^*f\|=\|g\|\leqslant\|T\|\cdot\|f\|\quad (\forall f\in\mathscr{Y}^*), \label{eq:::HIOUR32s2.5.25}
\end{align}
因此$T^*\in\mathscr{L}(\mathscr{Y}^*,\mathscr{X}^*)$，且
\begin{align*}
\|T^*\|\leqslant\|T\|. 
\end{align*}
注意到\eqref{eq:JHIODWI231122888}式左边由$f,T$完全确定,故$T$的共轭算子是唯一的.
\end{proof}

\begin{theorem}\label{theorem:泛函分析--等距同构基本定理}
映射$*:T\mapsto T^*$是$\mathscr{L}(\mathscr{X},\mathscr{Y})$到$\mathscr{L}(\mathscr{Y}^*,\mathscr{X}^*)$内的等距同构,进而$\|T^{*}\|=\|T\|$.
\end{theorem}
\begin{proof}
\begin{enumerate}[(1)]
\item 证$*:T\mapsto T^*$是线性的：
\begin{align*}
[(\alpha_1T_1+\alpha_2T_2)^*f](x)=f[(\alpha_1T_1+\alpha_2T_2)x]
=\alpha_1f(T_1x)+\alpha_2f(T_2x)=[(\alpha_1T_1^*+\alpha_2T_2^*)f](x)
\end{align*}
$(\forall x\in\mathscr{X},\forall f\in\mathscr{Y}^*,\forall\alpha_1,\alpha_2\in\mathbb{K})$.

\item 证等距。已有\eqref{eq:::HIOUR32s2.5.26}式，只需证$\|T\|\leqslant\|T^*\|$.

对$\forall x\in\mathscr{X}$，若$Tx\neq\theta$，由\refcor{corollary:泛函分析--推论2.4.6}，$\exists f\in\mathscr{Y}^*$，使得$f(Tx)=\|Tx\|$，且$\|f\|=1$. 于是由\refpro{proposition:有界线性算子的基本不等式}知
\begin{align*}
\|Tx\|=f(Tx)=(T^*f)(x)\leqslant\|T^*f\|\cdot\|x\|\leqslant\|T^*\|\cdot\|x\|,
\end{align*}
即$\|T\|\leqslant\|T^*\|$.
\end{enumerate}
\end{proof}

\begin{theorem}
设$\mathscr{X},\mathscr{Y}$是$B^*$空间，$T\in\mathscr{L}(\mathscr{X},\mathscr{Y})$，则$T^{**}\in\mathscr{L}(\mathscr{X}^{**},\mathscr{Y}^{**})$是$T$在$\mathscr{X}^{**}$上的延拓，且满足$\|T^{**}\|=\|T^{*}\|=\|T\|$.
\end{theorem}
\begin{proof}
$T^*$的共轭算子$T^{**}=(T^*)^*\in\mathscr{L}(\mathscr{X}^{**},\mathscr{Y}^{**})$. 由\refthe{theorem:原空间可等距嵌入第二共轭空间}知$\mathscr{X}\subseteq\mathscr{X}^{**},\mathscr{Y}\subseteq\mathscr{Y}^{**}$，设自然嵌入映射分别为$U$和$V$，则
\begin{align*}
\langle T^{**}Ux,f\rangle=\langle Ux,T^*f\rangle
=\langle T^*f,x\rangle=\langle f,Tx\rangle
=\langle VTx,f\rangle\quad (\forall f\in\mathscr{Y}^*,\forall x\in\mathscr{X}).
\end{align*}
故$T^{**}Ux=VTx$，即$T^{**}$是$T$在$\mathscr{X}^{**}$上的同构意义下的扩张.由\refthe{theorem:泛函分析--等距同构基本定理}知$\|T^{**}\|=\|T^*\|=\|T\|$.
\end{proof}

\begin{example}
设$(\Omega,\mathscr{B},\mu)$是一个测度空间，$K(x,y)$是$\Omega\times\Omega$上的二元平方可积函数：
\[
\iint_{\Omega\times\Omega}|K(x,y)|^2\mathrm{d}\mu(x)\mathrm{d}\mu(y)<\infty.
\]
定义算子
\[
T:u\mapsto (Tu)(x)=\int_{\Omega}K(x,y)u(y)\mathrm{d}\mu(y)\quad (\forall u\in L^2(\Omega,\mu)),
\]
则$T\in\mathscr{L}(L^2(\Omega,\mu))$，且
\[
(T^*v)(x)=\int_{\Omega}K(y,x)v(y)\mathrm{d}\mu(y)\quad (\forall v\in L^2(\Omega,\mu)).
\]
\end{example}

\begin{proof}
记$\|\cdot\|$为$L^2(\Omega,\mu)$上的范数，则
\begin{align*}
\|Tu\|^2&=\int_{\Omega}\left|\int_{\Omega}K(x,y)u(y)\mathrm{d}\mu(y)\right|^2\mathrm{d}\mu(x)\\
&\leqslant\int_{\Omega}\left(\int_{\Omega}|K(x,y)|^2\mathrm{d}\mu(y)\int_{\Omega}|u(y)|^2\mathrm{d}\mu(y)\right)\mathrm{d}\mu(x)\\
&\leqslant\left(\iint_{\Omega\times\Omega}|K(x,y)|^2\mathrm{d}\mu(x)\mathrm{d}\mu(y)\right)\|u\|^2\quad (\forall u\in L^2(\Omega,\mu)).
\end{align*}
且对$\forall u,v\in L^2(\Omega,\mu)$,有
\begin{align*}
\langle T^*v,u\rangle&=\langle v,Tu\rangle
=\int_{\Omega}\left(\int_{\Omega}K(x,y)u(y)\mathrm{d}\mu(y)\right)v(x)\mathrm{d}\mu(x)\\
&\stackrel{\one}{=}\iint_{\Omega\times\Omega}K(x,y)u(y)v(x)\mathrm{d}\mu(x)\mathrm{d}\mu(y)
=\int_{\Omega}\left(\int_{\Omega}K(x,y)v(x)\mathrm{d}\mu(x)\right)u(y)\mathrm{d}\mu(y)
\end{align*}
上式第三个等号成立是因
\[
\iint_{\Omega\times\Omega}|K(x,y)||u(y)||v(x)|\mathrm{d}\mu(x)\mathrm{d}\mu(y)\leqslant\|u\|\cdot\|v\|\left(\iint_{\Omega\times\Omega}|K(x,y)|^2\mathrm{d}\mu(x)\mathrm{d}\mu(y)\right)^{\frac{1}{2}}<\infty,
\]
故可应用Fubini定理.故
\[
(T^*v)(y)=\int_{\Omega}K(x,y)v(x)\mathrm{d}\mu(x)\quad (\forall v\in L^2(\Omega,\mu)).
\]
\end{proof}

\begin{lemma}[Young不等式]
设$f\in L^p(\mathbb{R})\ (1\leqslant p\leqslant\infty)$，$K\in L^1(\mathbb{R})$，则
\begin{align}
\|K*f\|_p\leqslant\|K\|_1\cdot\|f\|_p, \label{eq:::HIOUR32s2.5.27}
\end{align}
其中$\|\cdot\|_p$表示$L^p(\mathbb{R})$的范数$(1\leqslant p\leqslant\infty)$.
\end{lemma}
\begin{proof}
当$1<p<\infty$时，由Hölder不等式，
\begin{align*}
\left|\int_{-\infty}^{\infty}K(x-y)f(y)\mathrm{d}y\right|&\leqslant\int_{-\infty}^{\infty}|K(x-y)|^{\frac{1}{q}}\cdot|K(x-y)|^{\frac{1}{p}}|f(y)|\mathrm{d}y\\
&\leqslant\left(\int_{-\infty}^{\infty}|K(x-y)|\mathrm{d}y\right)^{\frac{1}{q}}\left(\int_{-\infty}^{\infty}|K(x-y)|\cdot|f(y)|^p\mathrm{d}y\right)^{\frac{1}{p}},
\end{align*}
而
\begin{align*}
\int_{-\infty}^{\infty}\left|\int_{-\infty}^{\infty}K(x-y)f(y)\mathrm{d}y\right|^p\mathrm{d}x\leqslant\|K\|_1^{\frac{p}{q}}\int_{-\infty}^{\infty}\int_{-\infty}^{\infty}|K(x-y)|\cdot|f(y)|^p\mathrm{d}y\mathrm{d}x=\|K\|_1^{\frac{p}{q}}\cdot\|f\|_p^p\cdot\|K\|_1.
\end{align*}
由Fubini定理，$(K*f)(x)$几乎处处存在有限，且\eqref{eq:::HIOUR32s2.5.27}式成立.

当$p=1$或$\infty$时，不等式\eqref{eq:::HIOUR32s2.5.27}显然成立.
\end{proof}

\begin{example}
设$K(x)$是$\mathbb{R}$上的$L^1$函数，考察空间$L^p(\mathbb{R})\ (1\leqslant p\leqslant\infty)$上的卷积算子
\[
(K*f)(x)\triangleq\int_{-\infty}^{\infty}K(x-y)f(y)\mathrm{d}y,
\]
求其共轭算子。
\end{example}
\begin{remark}
以上考虑的共轭空间和共轭算子都是在实数域上的共轭。若用复数域，则每个$L^q$函数$g$对应着$L^p$空间上的一个反连续线性泛函：
\[
F_g(f)=\int_{\Omega}f\cdot\overline{g}\mathrm{d}\mu.
\]
这时的复共轭算子为$T^*=\overline{\wideparen{K}}*$.
\end{remark}
\begin{proof}
记$\wideparen{K}(x)\triangleq K(-x)$，由Fubini定理，
\begin{align*}
\int_{-\infty}^{\infty}\left(\int_{-\infty}^{\infty}K(x-y)f(y)\mathrm{d}y\right)g(x)\mathrm{d}x=\int_{-\infty}^{\infty}f(y)\left(\int_{-\infty}^{\infty}K(x-y)g(x)\mathrm{d}x\right)\mathrm{d}y=\int_{-\infty}^{\infty}(\wideparen{K}*g)(y)f(y)\mathrm{d}y.
\end{align*}
故$T\triangleq K*$的共轭算子为$T^*=\wideparen{K}*$.
\end{proof}


\subsection{弱收敛及*弱收敛}

\begin{definition}
设$\mathscr{X}$是一个$B^*$空间，$\{x_n\}\subseteq\mathscr{X}$，$x\in\mathscr{X}$。称$\{x_n\}$\textbf{弱收敛}到$x$，记作$x_n\rightharpoonup x$，是指：对$\forall f\in\mathscr{X}^*$都有
\begin{align*}
\lim\limits_{n\to\infty} f(x_n)=f(x).
\end{align*}
这时$x$称作点列$\{x_n\}$的\textbf{弱极限}。

为区别起见，今后我们称$x_n\to x$(按范数收敛)为$\{x_n\}$\textbf{强收敛}到$x$，或$x$是$\{x_n\}$的\textbf{强极限}。
\end{definition}
\begin{remark}
若$\dim\mathscr{X}<\infty$，则弱收敛与强收敛是等价的。事实上，设$e_1,e_2,\cdots,e_m$是$\mathscr{X}$的一组基，并设
\begin{align*}
x_n&=\xi_1^{(n)}e_1+\xi_2^{(n)}e_2+\cdots+\xi_m^{(n)}e_m\quad(n=1,2,\cdots),\\
x&=\xi_1^{(0)}e_1+\xi_2^{(0)}e_2+\cdots+\xi_m^{(0)}e_m.
\end{align*}
取$f_i\in\mathscr{X}^*\ (i=1,2,\cdots,m)$，使得$f_i(e_j)=\delta_{ij}\ (i,j=1,2,\cdots,m)$，便有
$$f_i(x_n)=\xi_i^{(n)}\quad\text{与}\quad f_i(x)=\xi_i^{(0)}\quad(i=1,2,\cdots,m).$$
今若$x_n\rightharpoonup x$，则有$\lim\limits_{n\to\infty}f(x_n)=f(x)\ (\forall f\in\mathscr{X}^*)$，从而也有
\begin{align*}
\lim\limits_{n\to\infty} f_i(x_n)=f_i(x)\quad(i=1,2,\cdots,m),
\end{align*}
即
\begin{align*}
\lim\limits_{n\to\infty} \xi_i^{(n)}=\xi_i^{(0)}\quad(i=1,2,\cdots,m).
\end{align*}
换句话说就是$\{x_n\}$按其坐标收敛于$x$。反过来，若$\{x_n\}$按其坐标收敛于$x$，那么$x_n\to x\ (n\to\infty)$，由此便可推出$x_n\rightharpoonup x\ (n\to\infty)$。
\end{remark}

\begin{proposition}
\begin{enumerate}[(1)]
\item 弱极限若存在必唯一。
\item 强极限若存在必是弱极限。
\end{enumerate}
\end{proposition}
\begin{remark}
但反过来,当$\text{dim}\,\mathscr{X}=\infty$时,弱极限存在却未必有强极限.
\end{remark}
\begin{proof}
\begin{enumerate}[(1)]
\item 若有$x_n\rightharpoonup x,\ x_n\rightharpoonup y\ (n\to\infty)$，由定义推得
\begin{align*}
f(x)=\lim\limits_{n\to\infty}f(x_n)=f(y)\quad(\forall f\in\mathscr{X}^*).
\end{align*}
因此$f(x-y)=0$,利用\refcor{corollary:泛函分析--推论2.4.6}，即得$x-y=\theta$,于是$x=y$。

\item 若$x_n\to x\ (n\to\infty)$，则$\forall f\in\mathscr{X}^*$,由\refpro{proposition:有界线性算子的基本不等式}有
\begin{align*}
|f(x_n)-f(x)|\leqslant\|f\|\cdot\|x_n-x\|\to0\quad(n\to\infty),
\end{align*}
即得$\lim\limits_{n\to\infty}f(x_n)=f(x)$.故$x_n\rightharpoonup x\ (n\to\infty)$.
\end{enumerate}
\end{proof}

\begin{example}
在$L^2[0,1]$中，设$x_n=x_n(t)=\sin n\pi t$，则根据Riemann-Lebesgue引理，显然有
\begin{align*}
\langle f,x_n\rangle=\int_{0}^{1}f(t)\sin n\pi t\mathrm{d}t\to0\quad(\forall f\in L^2[0,1]),
\end{align*}
即$x_n\rightharpoonup\theta\ (n\to\infty)$。但$\|x_n\|=\frac{1}{\sqrt{2}}$，不可能有$x_n\to\theta\ (n\to\infty)$。
\end{example}

\begin{theorem}[Mazur定理]
设$\mathscr{X}$是一个$B^*$空间，$x_n\rightharpoonup x_0\ (n\to\infty)$，则$\forall\varepsilon>0,\ \exists\lambda_i\geqslant0\ (i=1,2,\cdots,n),\ \sum\limits_{i=1}^{n}\lambda_i=1$，使得
\begin{align*}
\left\|x_0-\sum\limits_{i=1}^{n}\lambda_ix_i\right\|\leqslant\varepsilon.
\end{align*}
\end{theorem}
\begin{remark}
这表明弱收敛确实与强收敛不同。然而反过来，若$x_n\rightharpoonup x\ (n\to\infty)$，我们却可以找到$\{x_n\}$的凸组合序列，使其强收敛到$x$。
\end{remark}
\begin{proof}
设$M\triangleq\overline{\mathrm{co}(\{x_n\})}$，则$M$是$\mathscr{X}$中的一个闭凸集。倘若$x_0\notin M$，应用\hyperref[corollary:泛函分析--Ascoli定理]{Ascoli定理}，$\exists f\in\mathscr{X}^*$及$\alpha\in\mathbb{R}$，使得
\begin{align*}
f(x)<\alpha<f(x_0)\quad(\forall x\in M).
\end{align*}
从而
\begin{align*}
f(x_n)<\alpha<f(x_0)\quad(\forall n\in\mathbb{N}).
\end{align*}
这与$x_n\rightharpoonup x_0\ (n\to\infty)$矛盾!
\end{proof}

\begin{definition}
设$\mathscr{X}$是$B^*$空间，$\{f_n\}\subseteq\mathscr{X}^*$，$f\in\mathscr{X}^*$。称$\{f_n\}$\textbf{*弱收敛}到$f$，记作$w^*-\lim\limits_{n\to\infty}f_n=f$或$f_n \xrightarrow{w^*} f$，是指：对于$\forall x\in\mathscr{X}$，都有$\lim\limits_{n\to\infty}f_n(x)=f(x)$。这时$f$称作泛函序列$\{f_n\}$的\textbf{*弱极限}。
\end{definition}

\begin{proposition}\label{proposition:共轭空间上的弱收敛可推出星弱收敛}
设 $\mathscr{X}$ 是$B^*$空间,
若 $\{f_n\} \subseteq \mathscr{X}^*$ 弱收敛于 $f \in \mathscr{X}^*$，则 $\{f_n\}$ 也 *弱收敛于 $f$。

特别地,若 $\mathscr{X}$ 是自反空间，则 $\mathscr{X}^*$ 中弱收敛与 *弱收敛等价。
\end{proposition}
\begin{proof}

由\refthe{theorem:原空间可等距嵌入第二共轭空间}可知,存在线性等距嵌入 $J: \mathscr{X} \hookrightarrow \mathscr{X}^{**}$，并且对每个 $x \in \mathscr{X}$,有
\begin{align*}
\left< Jx,f \right> =\left< f,x \right> ,\forall f\in\mathscr{X}^*\iff Jx\left( f \right) =f\left( x \right) ,\forall f\in \mathscr{X} ^*.
\end{align*}
因此可将 $\mathscr{X}$ 视为 $\mathscr{X}^{**}$ 的子空间.
 
设 $f_n \rightharpoonup f$，即对任意 $\Phi \in \mathscr{X}^{**}$ 有 $\Phi(f_n) \to \Phi(f)$。特别地，对每个 $x \in \mathscr{X}$，取 $\Phi = J(x) \in \mathscr{X}^{**}$，则
\begin{gather*}
\left< Jx,f_n \right> =Jx\left( f_n \right) \rightarrow Jx\left( f \right) =\left< Jx,f \right> \quad \left( n\rightarrow \infty \right) ;
\\
\left< Jx,f_n \right> =\left< f_n,x \right> =f_n\left( x \right) ,\quad \left< Jx,f \right> =\left< f,x \right> =f\left( x \right) .
\end{gather*}
故$f_n(x)\to f(x)(n\to \infty)$,对所有 $x\in\mathscr{X}$ 成立。这正是$w^*-\lim\limits_{n\to\infty}f_n=f$.

特别地,若$\mathscr{X}$是自反空间，即 $J(\mathscr{X}) = \mathscr{X}^{**}$。  
由上述证明已知弱收敛 $\Rightarrow$ *弱收敛.
反之，设 $f_n \xrightarrow{w^*} f$，即对每个 $x\in\mathscr{X}$ 有 $f_n(x) \to f(x)$。任取 $\Phi \in \mathscr{X}^{**}$，由自反性存在 $x_\Phi \in \mathscr{X}$ 使得 $\Phi = J(x_\Phi)$。于是
\[
\Phi(f_n) = J(x_\Phi)(f_n) = f_n(x_\Phi) \to f(x_\Phi) = J(x_\Phi)(f) = \Phi(f).
\]
因此 $f_n \rightharpoonup f$。故弱收敛与 *弱收敛等价.
\end{proof}

\begin{theorem}\label{theorem:弱收敛的充要条件111}
设$\mathscr{X}$是一个$B^*$空间，又设$\{x_n\}\subseteq\mathscr{X},\ x\in\mathscr{X}$，则为了$x_n\rightharpoonup x\ (n\to\infty)$当且仅当
\begin{enumerate}[(1)]
\item $\|x_n\|$有界；
\item 对$\mathscr{X}^*$中的一个稠密子集$M^*$上的一切$f$都有
\begin{align*}
\lim\limits_{n\to\infty}f(x_n)=f(x).
\end{align*}
\end{enumerate}
\end{theorem}
\begin{proof}
由\refthe{theorem:原空间可等距嵌入第二共轭空间}可知,只需把$x_n$看成是$\mathscr{X}^*$上的有界线性泛函：
$$\langle x_n,f\rangle\triangleq f(x_n)\quad(\forall f\in\mathscr{X}^*).$$
应用\hyperref[theorem:泛函分析--Banach-Steinhaus定理]{Banach-Steinhaus定理}即得结论.
\end{proof}

\begin{theorem}\label{theorem:泛函分析--定理2.5.21}
设$\mathscr{X}$是一个$B$空间，又设$\{f_n\}\subseteq\mathscr{X}^*,\ f\in\mathscr{X}^*$，则为了$w^*-\lim\limits_{n\to\infty}f_n=f$当且仅当
\begin{enumerate}[(1)]
\item $\|f_n\|$有界；
\item 对$\mathscr{X}$中的一个稠密子集$M$上的一切$x$都有
\begin{align*}
\lim\limits_{n\to\infty}f_n(x)=f(x).
\end{align*}
\end{enumerate}
\end{theorem}
\begin{proof}
本定理是\hyperref[theorem:泛函分析--Banach-Steinhaus定理]{Banach-Steinhaus定理}的特殊情形.
\end{proof}

\begin{proposition}
求证: 有限维$B^*$空间$X$必是自反的,且
\begin{align*}
\dim X=\dim X^{*}=\dim X^{**}<+\infty.
\end{align*}
进而$\mathscr{X}^*$ 中弱收敛与 *弱收敛等价.并且$X$中强收敛等价于弱收敛.
\end{proposition}
\begin{proof}
设 $\dim X=n$. 取 $X$ 的一组基 $e_1,\cdots,e_n$. 定义 $e_i^*\in X^*$ 为
\begin{align*}
e_i^*(e_j)=\delta_{ij}=\begin{cases}
1,&i=j,\\
0,&i\ne j\\
\end{cases},\quad i,j=1,\cdots,n.
\end{align*}
则对任意 $f\in X^*$, 有
\begin{align*}
f=\sum\limits_{i=1}^n f(e_i)e_i^*.
\end{align*}
事实上, 对任意 $x=\sum\limits_{j=1}^n x_j e_j\in X$, 有
\begin{align*}
\left(\sum\limits_{i=1}^n f(e_i)e_i^*\right)(x)
=\sum\limits_{i=1}^n f(e_i)x_i=f\left(\sum\limits_{i=1}^n x_i e_i\right)
=f(x).
\end{align*}
因此 $e_1^*,\cdots,e_n^*$ 张成 $X^*$. 又若 $\sum\limits_{i=1}^n a_i e_i^*=0$, 则对每个 $j$, 有
\begin{align*}
0=\left(\sum\limits_{i=1}^n a_i e_i^*\right)(e_j)=a_j.
\end{align*}
故 $e_1^*,\cdots,e_n^*$ 线性无关. 于是 $\dim X^*=n$. 同理可得 $\dim X^{**}=n$.

由\cref{theorem:原空间可等距嵌入第二共轭空间}，定义等距嵌入$J:X\to X^{**}$ 为
\begin{align*}
(Jx)(f)\triangleq f(x),\quad x\in X,\ f\in X^*.
\end{align*}
显然 $J$ 是线性映射. 若 $Jx=0$, 则对一切 $f\in X^*$ 有 $Jx(f)=f(x)=0$. 特别取 $f=e_i^*$, 得 $x_i=0$ 对每个 $i$ 成立, 故 $x=0$. 因此 $J$ 是单射. 由于 $\dim X=n=\dim X^{**}$, 故由\refpro{Basis of Algebra-proposition:像和原像空间维数相同时线性同构的充要条件}知线性单射 $J$ 必为满射, 从而 $J$ 是线性同构. 进而$J$ 是等距同构, 故 $X$ 自反.进而由\cref{proposition:共轭空间上的弱收敛可推出星弱收敛}知$\mathscr{X}^*$ 中弱收敛与 *弱收敛等价.

最后证明强收敛与弱收敛等价.
若 \(x_n\) 强收敛于 \(x\)，即 \(\|x_n-x\|\to 0\)，则对于任意的 \(f\in X^*\)，有
\begin{align*}
|f(x_n)-f(x)| = |f(x_n-x)| \leqslant \|f\| \|x_n-x\| \to 0.
\end{align*}
故 \(x_n\) 弱收敛于 \(x\).

反之，若 \(x_n\) 弱收敛于 \(x\)，即对于一切 \(f\in X^*\) 都有 \(f(x_n)\to f(x)\).
令 \(x_n = \sum\limits_{i=1}^n a_i^{(n)} e_i\)，\(x = \sum\limits_{i=1}^n a_i e_i\)。由于对偶基 \(e_i^*\in X^*\)，由弱收敛知 \(a_i^{(n)} = e_i^*(x_n) \to e_i^*(x) = a_i\) （对所有 \(i=1,\cdots,n\) 成立）.
于是
\begin{align*}
\|x_n-x\| = \left\|\sum\limits_{i=1}^n (a_i^{(n)} - a_i) e_i\right\| \leqslant \sum\limits_{i=1}^n |a_i^{(n)} - a_i| \|e_i\| \to 0.
\end{align*}
这说明 \(x_n\) 强收敛于 \(x\). 综合上述，在有限维空间中强收敛等价于弱收敛.
\end{proof}

\begin{proposition}
设 \(H\) 是一个Hilbert空间，\(\{x_n\}_{n=1}^{\infty} \subseteq H\) 是 \(H\) 中的点列，且 \(x \in H\)。则 $x_n\to x\ (n\to\infty)$的充要条件是
\begin{enumerate}[(1)]
\item $\|x_n\|\to\|x\|\ (n\to\infty)$;
\item $x_n\rightharpoonup x\ (n\to\infty)$.
\end{enumerate}
\end{proposition}
\begin{proof}

{\heiti 必要性:} 若 \(x_n\) 强收敛于 \(x\)，即 \(\|x_n - x\| \to 0\)。
先证范数收敛。由\hyperref[proposition:内积的连续性]{内积的连续性}可知
\begin{align*}
|\|x_n\| - \|x\|| \leqslant \|x_n - x\| \to 0.
\end{align*}
故 \(\|x_n\| \to \|x\|\)。
再证弱收敛。由\hyperref[theorem:Riesz表示定理-泛函]{Riesz表示定理}可知，弱收敛等价于对任意固定的 \(y \in H\)，有 \((x_n, y) \to (x, y)\)。因此，对任意固定的 \(y \in H\)，由Cauchy-Schwarz不等式，有
\begin{align}
|(x_n, y) - (x, y)| \leqslant \|x_n - x\| \cdot \|y\| \to 0.
\end{align}
因此 \((x_n, y) \to (x, y)\)，即 \(x_n\) 弱收敛于 \(x\)。

{\heiti 充分性:} 若 \(x_n\) 弱收敛于 \(x\)，且 \(\|x_n\| \to \|x\|\)。
由\hyperref[theorem:Riesz表示定理-泛函]{Riesz表示定理}可知，弱收敛等价于对任意固定的 \(y \in H\)，有 \((x_n, y) \to (x, y)\)。特别地，取 \(y = x\)，则有
\begin{align}
\label{eq:weak_conv_special}
(x_n, x) \to (x, x) = \|x\|^{2}.
\end{align}
由\hyperref[proposition:内积的连续性]{内积的连续性}可知
\begin{align}
\label{eq:norm_expand}
\|x_n - x\|^2 = \|x_n\|^2 - 2\text{Re}(x_n, x) + \|x\|^{2}.
\end{align}
由于已知 \(\|x_n\| \to \|x\|\)，即 \(\|x_n\|^2 \to \|x\|^2\)，再结合 \eqref{eq:weak_conv_special} 中的内积收敛性，对等式 \eqref{eq:norm_expand} 两边取极限，可得：
\begin{align*}
\lim_{n\to\infty} \|x_n - x\|^2 = \|x\|^2 - 2\|x\|^2 + \|x\|^2 = 0.
\end{align*}
因此 \(\lim_{n\to\infty} \|x_n - x\| = 0\)，即 \(x_n\) 强收敛于 \(x\).
\end{proof}

\begin{example}
在 \(L^2[0,1]\) 空间中考察点列 \(\{x_n(t) = \sqrt{2}\sin(n\pi t)\}\)。由黎曼-勒贝格引理可知，对于任意 \(f \in L^2[0,1]\)，都有
\begin{align*}
\lim_{n\to\infty} \int_{0}^{1} f(t)\sqrt{2}\sin(n\pi t) \mathrm{d}t = 0.
\end{align*}
因此，\(x_n \rightharpoonup \theta\)（弱收敛于零）。
然而，由于 \(\|x_n\|_{L^2}^2 = 2\int_0^1 \sin^2(n\pi t) \mathrm{d}t = 1\)，故 \(\|x_n\|\) 并不收敛于 \(\|\theta\| = 0\)。
因此，虽然 \(x_n\) 弱收敛于零，但并不强收敛于零。这说明命题中“\(\|x_n\| \to \|x\|\)”这一条件是不可或缺的。
\end{example}


\begin{definition}
设$\mathscr{X},\mathscr{Y}$是$B^*$空间。又设$T_n\ (n=1,2,\cdots),\ T\in\mathscr{L}(\mathscr{X},\mathscr{Y})$。
\begin{enumerate}[(1)]
\item 若$\|T_n-T\|\to0$，则称$T_n$\textbf{一致收敛}于$T$，记作$T_n\rightrightarrows T$。这时$T$称作$\{T_n\}$的\textbf{一致极限}。

\item 若$\|(T_n-T)x\|\to0\ (\forall x\in\mathscr{X})$，则称$T_n$\textbf{强收敛}于$T$，记作$T_n\to T$。这时$T$称作$\{T_n\}$的\textbf{强极限}。

\item 如果对于$\forall x\in\mathscr{X}$，以及$\forall f\in\mathscr{Y}^*$都有
\begin{align*}
\lim\limits_{n\to\infty}f(T_nx)=f(Tx),
\end{align*}
则称$T_n$\textbf{弱收敛}于$T$，记作$T_n\rightharpoonup T$。这时$T$称作$\{T_n\}$的\textbf{弱极限}。
\end{enumerate}
\end{definition}
\begin{remark}
显然，一致收敛$\implies$强收敛$\implies$弱收敛，而且每种极限若存在必是唯一的。但反过来一般不对。
\end{remark}

\begin{example}[强收敛而不一致收敛]
 
在空间$l^2$上考察左推移算子
$$T:x=(x_1,x_2,\cdots,x_n,\cdots)\mapsto Tx=(x_2,x_3,\cdots,x_n,\cdots).$$
令$T_n\triangleq T^n$，便有
$$T_nx=(x_{n+1},x_{n+2},\cdots)\quad(\forall x=(x_1,x_2,\cdots,x_n,\cdots)\in l^2).$$
可证：$T_n\to0$，但$T_n\to 0\ (n\to\infty)$。事实上，若取
$$e_n=(\underbrace{0,0,\cdots,0}_{n},1,0,\cdots),$$
那么$T_ne_{n+1}=e_1$，并且$\|e_n\|=1\ (\forall n\in\mathbb{N})$。因此
$$\|T_n\|\geqslant\|T_n(e_{n+1})\|=1,$$
从而$T_n\to 0$。但是对$\forall x=(x_1,x_2,\cdots,x_n,\cdots)\in l^2$有
\begin{align*}
\|T_nx\|=\left(\sum\limits_{i=1}^{\infty}|x_{n+i}|^2\right)^{\frac{1}{2}}\to0\quad(n\to\infty),
\end{align*}
即$T_n\to0$。
\end{example}

\begin{example}[弱收敛而不强收敛]

在空间$l^2$上考察右推移算子
$$S:x=(x_1,x_2,\cdots,x_n,\cdots)\mapsto Sx=(0,x_1,x_2,\cdots,x_n,\cdots).$$
令$S_n\triangleq S^n$，便有
$$S_nx=(\underbrace{0,0,\cdots,0}_{n},x_1,x_2,\cdots)\quad(\forall x\in l^2).$$
显然，$\|S_nx\|=\|x\|\ (\forall x\in l^2)$，从而$S_n\nrightarrow0$。但是对于$\forall f=(y_1,y_2,\cdots,y_n,\cdots)\in(l^2)^*=l^2$，我们有
\begin{align*}
|\langle f,S_nx\rangle|=\left|\sum\limits_{i=1}^{\infty}y_{i+n}x_i\right|\leqslant\left(\sum\limits_{i=1}^{\infty}|y_{i+n}|^2\right)^{\frac{1}{2}}\|x\|\to0\quad(n\to\infty),
\end{align*}
即$S_n\rightharpoonup0\ (n\to\infty)$。
\end{example}

\subsection{弱列紧性与*弱列紧性}

\begin{definition}
设$\mathscr{X}$是$B^*$空间,$A\subset \mathscr{X},A^*\subset \mathscr{X}^*$,称集$A$是\textbf{弱列紧的}，是指$A$中的任意点列有一个弱收敛子列。
又称$A^*$是*\textbf{弱列紧的}，是指$A^*$中的任意点列有一个*弱收敛的子列.
\end{definition}

\begin{theorem}
设$\mathscr{X}$是可分的$B^*$空间，那么$\mathscr{X}^*$上的任意有界列$\{f_n\}$必有*弱收敛的子列。
\end{theorem}
\begin{proof}
因为$\mathscr{X}$可分，所以$\mathscr{X}$有可数的稠密子集$\{x_m\}$。因为$\{f_n\}$有界，所以对每一个固定的$m$，数集
$$\{\langle f_n,x_m\rangle|n,m\in\mathbb{N}\}\subseteq \mathbb{K}=\mathbb{C}\text{或}\mathbb{R}$$
是有界的。由\hyperref[Complex Analysis-theorem:Bolzano-Weierstrass定理]{Bolzano-Weierstrass定理}和\hyperref[proposition:对角线法则]{对角线法则},存在子列$\{f_{n_k}\}_{k=1}^{\infty}$，使得对$\forall m\in\mathbb{N}$，
$$\{\langle f_{n_k},x_m\rangle\}_{k=1}^{\infty}$$
是收敛数列。再由$\{x_m\}$在$\mathscr{X}$中稠密，以及$\{f_n\}$有界，可见对于$\forall x\in\mathscr{X}$，
$$\{\langle f_{n_k},x\rangle\}_{k=1}^{\infty}$$
是收敛数列。记$f(x)\triangleq\lim\limits_{k\to\infty}\langle f_{n_k},x\rangle$。易见$f$是线性的，并且由\refpro{proposition:有界线性算子的基本不等式}知
$$|f(x)|\leqslant\sup\limits_{n}\|f_n\|\cdot\|x\|\quad(\forall x\in\mathscr{X}).$$
从而有$f\in\mathscr{X}^*$，使得
$$f(x)=\lim\limits_{k\to\infty}\langle f_{n_k},x\rangle\quad(\forall x\in\mathscr{X}).$$
即得$w^*-\lim\limits_{k\to\infty}f_{n_k}=f$。
\end{proof}

\begin{theorem}[Banach定理]\label{theorem:泛函分析------Banach定理}
设$\mathscr{X}$是$B^*$空间。若$\mathscr{X}$的共轭空间$\mathscr{X}^*$是可分的，则$\mathscr{X}$本身必是可分的。
\end{theorem}
\begin{proof}
\begin{enumerate}[(1)]
\item 考察$\mathscr{X}^*$的单位球面$S_1^*\triangleq\{f\in\mathscr{X}^*|\|f\|=1\}$。$S_1^*$是可分的。事实上，由$\mathscr{X}^*$可分，$\exists\{f_n\}\subseteq\mathscr{X}^*$，使得$\forall f\in S_1^*,\ \exists\{n_k\}_{k=1}^{\infty}$，使得
\begin{align*}
\lim\limits_{k\to\infty}f_{n_k}=f.
\end{align*}
今令$g_n\triangleq f_n/\|f_n\|$（不妨设$f_n\neq\theta$），便有
\begin{align*}
\|f-g_{n_k}\|\leqslant\|f-f_{n_k}\|+\|f_{n_k}-g_{n_k}\|=\|f-f_{n_k}\|+|1-\|f_{n_k}\||\to0\quad(k\to\infty).
\end{align*}
由此可见$S_1^*$有可数的稠密子集$\{g_n\}$。

\item 对于每个$g_n$，因为$\|g_n\|=1$，所以可以选取$x_n\in\mathscr{X}$，使得
$$\|x_n\|=1,\quad\text{并且}\quad g_n(x_n)\geqslant\frac{1}{2}.$$
记$\mathscr{X}_0\triangleq\overline{\mathrm{span}\{x_n\}}$，它显然是可分的（$x_n$的有理系数的线性组合在$\mathscr{X}_0$中稠密）。

\item 证明$\mathscr{X}_0=\mathscr{X}$。倘若不然，存在$x_0\in\mathscr{X}\setminus\mathscr{X}_0$，不妨设$\|x_0\|=1$，利用\refthe{theorem:泛函分析--零值泛函延拓定理}，$\exists f_0\in\mathscr{X}^*$，使得$\|f_0\|=1$，从而$f_0\in S_1^*$，并且$f_0(x)=0\ (\forall x\in\mathscr{X}_0)$。
对此$f_0$我们有
\begin{align*}
\|g_n-f_0\|=\sup\limits_{\|x\|=1}|g_n(x)-f_0(x)|
\geqslant|g_n(x_n)-f_0(x_n)|=|g_n(x_n)|\geqslant\frac{1}{2}.
\end{align*}
这与$\{g_n\}$在$S_1^*$中稠密相矛盾，即得$\mathscr{X}=\mathscr{X}_0$。从而$\mathscr{X}$是可分的。
\end{enumerate}
\end{proof}

\begin{theorem}[Pettis定理]\label{theorem:Pettis定理}
自反空间$\mathscr{X}$的闭子空间$\mathscr{X}_0$必是自反空间。
\end{theorem}
\begin{proof}
要证：若$z_0\in\mathscr{X}_0^{**}$，则必$z_0\in\mathscr{X}_0$；也就是要证：$\exists x\in\mathscr{X}_0$，使得
\begin{align}
\langle z_0,f_0\rangle=\langle f_0,x\rangle\quad(\forall f_0\in\mathscr{X}_0^*).\label{eq:::HIOUR32s2.5.28}
\end{align}
今对$\forall f\in\mathscr{X}^*$，考察$f$在$\mathscr{X}_0$上的限制$Tf=f_0\in\mathscr{X}_0^*$。因为
$$\|f_0\|\leqslant\|f\|,$$
所以$T\in\mathscr{L}(\mathscr{X}^*,\mathscr{X}_0^*)$。于是$z\triangleq T^*z_0\in\mathscr{X}^{**}$，又$\mathscr{X}$自反，因此$\exists x\in\mathscr{X}$，使得
\begin{align}
\langle z,f\rangle=\langle f,x\rangle\quad(\forall f\in\mathscr{X}^*).\label{eq:::HIOUR32s2.5.29}
\end{align}
今证此$x\in\mathscr{X}_0$。倘若不然，由\refthe{theorem:泛函分析--零值泛函延拓定理}，$\exists f\in\mathscr{X}^*$，使得
$$f(\mathscr{X}_0)=0,\quad\text{且}\quad\langle f,x\rangle=1,$$
从而$Tf=\theta$。但这导出矛盾：
$$0=\langle z_0,Tf\rangle=\langle T^*z_0,f\rangle=\langle z,f\rangle=\langle f,x\rangle=1.$$
这就证明了$\exists x\in\mathscr{X}_0$，使得\eqref{eq:::HIOUR32s2.5.29}式成立。现在要证此$x$还适合\eqref{eq:::HIOUR32s2.5.28}式。事实上，$\forall f_0\in\mathscr{X}_0^*$，由\hyperref[theorem:泛函分析--(哈恩-巴拿赫)Hahn-Banach定理]{Hahn-Banach定理}，存在$f\in\mathscr{X}^*$，使得$f_0=Tf$。从而我们有
$$\langle z_0,f_0\rangle=\langle z_0,Tf\rangle=\langle z,f\rangle,$$
以及
$$\langle f_0,x\rangle=\langle f,x\rangle\quad(\forall x\in\mathscr{X}_0).$$
由此可见，适合\eqref{eq:::HIOUR32s2.5.29}式的$x$必适合\eqref{eq:::HIOUR32s2.5.28}式。
\end{proof}

\begin{theorem}[Eberlein-Smulian定理]\label{theorem:Eberlein-Smulian定理}
自反空间的单位(闭)球是弱(自)列紧的。
\end{theorem}
\begin{proof}
\begin{enumerate}[(1)]
\item 先证：自反空间$\mathscr{X}$中的任何有界点列$\{x_n\}$必有一个在$\mathscr{X}$中弱收敛的子列。令
$$\mathscr{X}_0\triangleq\overline{\mathrm{span}\{x_n\}}.$$
因为$\mathscr{X}$自反，所以$\mathscr{X}_0$也是自反的。又显然$\mathscr{X}_0$是可分的，这表明$(\mathscr{X}_0^*)^*$是可分的。再由\hyperref[theorem:泛函分析------Banach定理]{Banach定理}，$\mathscr{X}_0^*$也是可分的。若记$\{g_n\}$为$\mathscr{X}_0^{**}$中的元素，它适合
\begin{align}
\langle g_n,f\rangle=\langle f,x_n\rangle\quad(\forall f\in\mathscr{X}_0^*),\label{eq:::HIOUR32s2.5.30}
\end{align}
则$\{\|g_n\|\}$有界。设$M_0^*$是$\mathscr{X}_0^*$中的可数稠密子集，用\hyperref[proposition:对角线法则]{对角线法则}，在$\{g_n\}$中可以抽出一子列$\{g_{n_k}\}$及$\exists g\in\mathscr{X}_0^{**}$，使得
\begin{align}
\lim\limits_{k\to\infty}\langle g_{n_k},f\rangle=\langle g,f\rangle\quad(\forall f\in M_0^*).\label{eq:::HIOUR32s2.5.31}
\end{align}
再依\refthe{theorem:泛函分析--定理2.5.21}，\eqref{eq:::HIOUR32s2.5.31}式蕴含
\begin{align}
\lim\limits_{k\to\infty}\langle g_{n_k},f\rangle=\langle g,f\rangle\quad(\forall f\in\mathscr{X}_0^*).\label{eq:::HIOUR32s2.5.32}
\end{align}
还由$\mathscr{X}_0$的自反性，$\exists x_0\in\mathscr{X}_0$适合
\begin{align}
\langle g,f\rangle=\langle f,x_0\rangle\quad(\forall f\in\mathscr{X}_0^*).\label{eq:::HIOUR32s2.5.33}
\end{align}
联合\eqref{eq:::HIOUR32s2.5.30}式，\eqref{eq:::HIOUR32s2.5.32}式与\eqref{eq:::HIOUR32s2.5.33}式，便得到
\begin{align*}
\lim\limits_{k\to\infty}\langle f,x_{n_k}\rangle&=\lim\limits_{k\to\infty}\langle g_{n_k},f\rangle\\
&=\langle f,x_0\rangle\quad(\forall f\in\mathscr{X}_0^*).
\end{align*}
进一步对$\forall\widetilde{f}\in\mathscr{X}^*$，记$f\triangleq T\widetilde{f}$为$\widetilde{f}$在$\mathscr{X}_0$上的限制。因为$\{x_{n_k}\}\subseteq\mathscr{X}_0$及$x_0\in\mathscr{X}_0$，依上式便有
\begin{align*}
\lim\limits_{k\to\infty}\langle\widetilde{f},x_{n_k}\rangle=\lim\limits_{k\to\infty}\langle f,x_{n_k}\rangle
=\langle f,x_0\rangle=\langle\widetilde{f},x_0\rangle\quad(\forall\widetilde{f}\in\mathscr{X}^*),
\end{align*}
即得$x_{n_k}\rightharpoonup x_0$。于是$\mathscr{X}$中的任意有界集是弱列紧集，特别是单位球是弱列紧的。同样，单位闭球也是弱列紧的。

\item 证明单位闭球是弱自列紧的。设$x_{n_k}\rightharpoonup x_0$，并且$\|x_{n_k}\|\leqslant1$。由\refcor{corollary:泛函分析--推论2.4.6}，$\exists f\in\mathscr{X}^*$适合
$$f(x_0)=\|x_0\|,\quad\text{且}\quad\|f\|=1.$$
因此，我们有
\begin{align*}
\|x_0\|=f(x_0)=\lim\limits_{k\to\infty}f(x_{n_k})\leqslant\|f\|\sup\limits_{k\geqslant1}\|x_{n_k}\|\leqslant1,
\end{align*}
即$x_0$也在单位闭球内，从而单位闭球是弱自列紧的。
\end{enumerate}
\end{proof}

\begin{theorem}[Alaoglu定理]
设$\mathscr{X}$是$B^*$空间，则$\mathscr{X}^*$中的单位闭球是$*$弱紧的。
\end{theorem}
\begin{remark}
$*$弱紧的定义及证明可参见下册第五章。
\end{remark}

\begin{definition}
设$f\in L^1[0,2\pi]$，称
\begin{align*}
c_n\triangleq\frac{1}{2\pi}\int_{0}^{2\pi}f(x)\mathrm{e}^{-\mathrm{i}nx}\mathrm{d}x\quad(n=0,\pm1,\pm2,\cdots)
\end{align*}
为$f$的$\mathbf{Fourier}$\textbf{系数}；称级数
\begin{align}
\sum\limits_{n=-\infty}^{\infty}c_n\mathrm{e}^{\mathrm{i}nx}\label{eq:fourier_series}
\end{align}
为$f$的$\mathbf{Fourier}$\textbf{级数}。
\end{definition}

\begin{definition}
Fourier级数\eqref{eq:fourier_series}的$\mathbf{Cesaro}$\textbf{部分和（算术平均和）}定义为
\begin{align}
\sigma_n(f)(x)=\frac{1}{n+1}\sum\limits_{k=0}^{n}S_k(f)(x)
=\sum\limits_{k=-n}^{n}c_k\left(1-\frac{|k|}{n+1}\right)\mathrm{e}^{\mathrm{i}kx},\label{eq:cesaro_sum}
\end{align}
其中$S_n(f)(x)\triangleq\sum\limits_{k=-n}^{n}c_k\mathrm{e}^{\mathrm{i}kx}$。
\end{definition}

\begin{definition}[Fejer核]
$\mathbf{Fejer}$\textbf{核}定义为
\begin{align*}
K_n(x)\triangleq\frac{1}{2\pi}\sum\limits_{k=-n}^{n}\left(1-\frac{|k|}{n+1}\right)\mathrm{e}^{\mathrm{i}kx}
=\frac{1}{2\pi(n+1)}\left(\frac{\sin\frac{(n+1)x}{2}}{\sin\frac{x}{2}}\right)^2,
\end{align*}
满足卷积关系$\sigma_n(f)(x)=\int_{0}^{2\pi}f(y)K_n(x-y)\mathrm{d}y$，且$\int_{0}^{2\pi}K_n(x)\mathrm{d}x=1$。
\end{definition}
\begin{remark}
由$K_n$的非负性及Young不等式，对任意$f\in L^p[0,2\pi](1\leqslant p<\infty)$，成立$\|\sigma_n(f)\|_{L^p}\leqslant\|f\|_{L^p}$，即Cesaro部分和的$L^p$范数一致有界。
\end{remark}

\begin{theorem}
若$1<p\leqslant\infty$，且Fourier级数的Cesaro部分和满足
\begin{align}
\sup\limits_{n\geqslant1}\|\sigma_n\|_{L^p}<\infty,\label{eq:unif_bound}
\end{align}
则必存在$f\in L^p[0,2\pi]$，使得$\sigma_n$是$f$的Fourier级数的Cesaro部分和。
\end{theorem}
\begin{proof}
由共轭空间结论$L^p[0,2\pi]=L^q[0,2\pi]^*$（$\frac{1}{p}+\frac{1}{q}=1$），且$L^q[0,2\pi]$可分。结合条件\eqref{eq:unif_bound}，存在$f\in L^p[0,2\pi]$及子列$\{n_k\}$，使得
\begin{align*}
w^*-\lim\limits_{k\to\infty}\sigma_{n_k}(x)=f(x)\quad(\text{在 }L^p[0,2\pi]\text{ 中}).
\end{align*}
因$\mathrm{e}^{\mathrm{i}mx}\in L^q[0,2\pi]\ (m=0,\pm1,\pm2,\cdots)$，故
\begin{align*}
\frac{1}{2\pi}\int_{0}^{2\pi}f(x)\mathrm{e}^{-\mathrm{i}mx}\mathrm{d}x=\lim\limits_{k\to\infty}\frac{1}{2\pi}\int_{0}^{2\pi}\sigma_{n_k}(x)\mathrm{e}^{-\mathrm{i}mx}\mathrm{d}x
=\lim\limits_{k\to\infty}\left(1-\frac{|m|}{n_k+1}\right)c_m=c_m.
\end{align*}
因此$\sigma_n(x)$是$f$的Fourier级数的Cesaro部分和$\sigma_n(f)(x)$。
\end{proof}

\subsection{弱收敛的例子}

\begin{example}[振荡型弱收敛]

\begin{enumerate}[(1)]
\item 令$u_n(x)=\sin n\pi x$，由Riemann-Lebesgue引理，当$n\to\infty$时，$u_n$在$L^2[0,1]$中弱收敛但不强收敛于$0$。
\item 锯齿型函数列：
\begin{align*}
u_n(x)=
\begin{cases}
x-\frac{k}{n},& x\in\left[\frac{k}{n},\frac{2k+1}{2n}\right],\\
-x+\frac{k+1}{n},& x\in\left[\frac{2k+1}{2n},\frac{k+1}{n}\right],
\end{cases}
\end{align*}
该序列在$H^{1,2}(0,1)$中弱收敛但不强收敛于$0$，根源为函数列在$(0,1)$内剧烈振荡，常作为变分学中泛函的不收敛极小化子列。
\end{enumerate}
\begin{figure}[H]
\centering
\includegraphics[scale=0.3]{锯齿形极小化序列.png}
\caption{锯齿形极小化序列$u_n$}
\label{figure:zigzag_min_seq}
\end{figure}
该序列自身及逐段导数均有界。
\end{example}

\begin{example}[平移型弱收敛]

设$f\in L^p(\mathbb{R})\setminus\{0\},\ 1<p<\infty$，令$f_n(x)=f(x+n),\ x\in\mathbb{R},\ n=1,2,\cdots$，则$f_n\rightharpoonup0$，但$\|f_n\|=\|f\|$，即$f_n$弱收敛但不强收敛于$0$。
\end{example}
\begin{proof}
需证对任意$g\in(L^p(\mathbb{R}))^*=L^q(\mathbb{R})$（$\frac{1}{p}+\frac{1}{q}=1$），成立
\begin{align}
\int_{\mathbb{R}}f_n(x)g(x)\mathrm{d}x=\int_{\mathbb{R}}f(x+n)g(x)\mathrm{d}x\to0,\quad n\to\infty.\label{eq:trans_weak_conv}
\end{align}
$C_0^\infty(\mathbb{R})$在$L^q(\mathbb{R})$中稠密，由Banach-Steinhaus定理，只需对$g\in C_0^\infty(\mathbb{R})$验证\eqref{eq:trans_weak_conv}。作变量代换得
\begin{align*}
\int_{\mathbb{R}}f(x+n)g(x)\mathrm{d}x=\int_{\mathbb{R}}f(x)g(x-n)\mathrm{d}x.
\end{align*}

先设$f\in C_0^\infty(\mathbb{R})$，令$\varphi_n(x)=f(x)g(x-n)$，则存在常数$C_g>0$，使得$|\varphi_n(x)|\leqslant C_g|f(x)|$，且$\varphi_n(x)\to0,\text{a.e.},\ n\to\infty$。由Lebesgue控制收敛定理
\begin{align}
\lim\limits_{n\to\infty}\int_{\mathbb{R}}f_n(x)g(x)\mathrm{d}x=\lim\limits_{n\to\infty}\int_{\mathbb{R}}f(x)g(x-n)\mathrm{d}x=0.\label{eq:trans_conv_lim}
\end{align}

对一般$f$，$\forall\varepsilon>0$，取$f_\varepsilon\in C_0^\infty(\mathbb{R})$，满足$\|f-f_\varepsilon\|_{L^p}<\frac{\varepsilon}{2(\|g\|_{L^q}+1)}$，则
\begin{align*}
\left|\int_{\mathbb{R}}f_n(x)g(x)\mathrm{d}x\right|&=\left|\int_{\mathbb{R}}f(x)g(x-n)\mathrm{d}x\right|\\
&\leqslant\|f-f_\varepsilon\|_{L^p}\cdot\|g\|_{L^q}+\left|\int_{\mathbb{R}}f_\varepsilon(x)g(x-n)\mathrm{d}x\right|\\
&<\frac{\varepsilon}{2}+\left|\int_{\mathbb{R}}f_\varepsilon(x)g(x-n)\mathrm{d}x\right|.
\end{align*}
结合\eqref{eq:trans_conv_lim}，存在$N$，当$n\geqslant N$时，$\left|\int_{\mathbb{R}}f_n(x)g(x)\mathrm{d}x\right|<\varepsilon$，故\eqref{eq:trans_weak_conv}成立。
\end{proof}


\begin{example}[集中型弱收敛]

定义平面旋转不变函数空间$\mathscr{X}$：元素为$f(x,y)=u(r)$（$r=\sqrt{x^2+y^2}$），满足$u\in L^2_{\frac{4r\mathrm{d}r}{(1+r^2)^2}}(\mathbb{R}_+^1)$，$u'\in L^2_{r\mathrm{d}r}(\mathbb{R}_+^1)$，范数平方为
\begin{align*}
\|u\|^2=\int_{0}^{\infty}|u(r)|^2\frac{4r\mathrm{d}r}{(1+r^2)^2}+\int_{0}^{\infty}|u'(r)|^2r\mathrm{d}r.
\end{align*}

对$\lambda>0$，定义
\begin{align*}
\phi_\lambda(r)=\frac{-1+\lambda^2r^2}{1+\lambda^2r^2},\quad \psi_\lambda(r)=\frac{2\lambda r}{1+\lambda^2r^2},
\end{align*}
可验证$(\phi_\lambda,\psi_\lambda)\in\mathscr{X}$：
\begin{align*}
|\phi_\lambda(r)|^2+|\psi_\lambda(r)|^2=1\implies\int_{0}^{\infty}\left(|\phi_\lambda(r)|^2+|\psi_\lambda(r)|^2\right)\frac{4r\mathrm{d}r}{(1+r^2)^2}=2,
\end{align*}
且
\begin{align*}
\begin{cases}
(\phi_\lambda')_r=\frac{4\lambda^2r}{(1+\lambda^2r^2)^2},\\
(\psi_\lambda')_r=\frac{2\lambda(1-\lambda^2r^2)}{(1+\lambda^2r^2)^2},
\end{cases}
\end{align*}
故
\begin{align*}
\int_{0}^{\infty}\left(|(\phi_\lambda')_r|^2+|(\psi_\lambda')_r|^2\right)r\mathrm{d}r=2.
\end{align*}

当$\lambda_j\to\infty$时，$(\phi_{\lambda_j},\psi_{\lambda_j})$存在弱收敛子列，但$\phi_{\lambda_j}$不强收敛；除$r=0$外，$(\phi_\lambda,\psi_\lambda)$逐点收敛到常值函数$(1,0)$。该例子的几何意义为：一族能量有界的调和映射，能量可集中于一点。
\end{example}

\begin{example}
设$C$是收敛数列的全体，赋以范数
\begin{align*}
\|\cdot\|:\{\xi_k\}\in C\mapsto \sup\limits_{k\geqslant1}|\xi_k|,
\end{align*}
求证: $C^*=l^1$.
\end{example}

\begin{example}
设$C_0$是以$0$为极限的数列全体，赋以范数
\begin{align*}
\|\cdot\|:\{\xi_k\}\in C\mapsto \sup\limits_{k\geqslant1}|\xi_k|,
\end{align*}
求证: $C_0^*=l^1$.
\end{example}

\begin{proposition}\label{proposition:值域闭等价完备}
设$\mathscr{X}$是$B^*$空间，$T$是从$\mathscr{X}$到$\mathscr{X}^{**}$的自然映射，则 $R(T)$是闭的充要条件是$\mathscr{X}$是完备的.
\end{proposition}
\begin{proof}

由\cref{theorem:原空间可等距嵌入第二共轭空间}知$T$是从$\mathscr{X}$到$R(T)$的等距同构.

{\heiti 必要性:}由\subcref{proposition:完备度量空间的子空间等价于闭集}{proposition:完备度量空间的子空间等价于闭集-1}知$R(T)$是$\mathscr{X}^**$的完备度量子空间。又因为$\mathscr{X}$与$R(T)$等距同构，所以$\mathscr{X}$也是完备的。

{\heiti 充分性:}由\subcref{proposition:完备度量空间的子空间等价于闭集}{proposition:完备度量空间的子空间等价于闭集-2}知$\mathscr{X}$是闭的。又因为$\mathscr{X}$与$R(T)$等距同构，所以$R(T)$也是闭的。
\end{proof}

\begin{proposition}
求证: $B$空间是自反的当且仅当它的共轭空间是自反的.
\end{proposition}
\begin{proof}
设 $X$ 是 Banach 空间. 由\cref{theorem:原空间可等距嵌入第二共轭空间}，定义等距嵌入
\begin{gather*}
J_X:X\to X^{**},\quad (J_X x)(f)=f(x),\quad x\in X,\ f\in X^*,\\
J_{X^*}:X^*\to X^{***},\quad (J_{X^*}(f))(x^{**})=x^{**}(f),\quad f\in X^*,\ x^{**}\in X^{**}.
\end{gather*}
因此只需证明$J_X$是满射等价于$J_{X^*}$是满射.

{\heiti 必要性:} 假设 $X$ 自反, 即 $J_X$ 是满射. 任取 $x^{***}\in X^{***}$. 定义 $f=x^{***}\circ J_X$, 则 $f\in X^*$. 对任意 $x^{**}\in X^{**}$, 由 $J_X$ 满射, 存在 $x\in X$ 使得 $x^{**}=J_Xx$. 于是
\begin{align*}
J_{X^*}(f)(x^{**})=x^{**}(f)=(J_Xx)(f)=f(x)=x^{***}(J_Xx)=x^{***}(x^{**}).
\end{align*}
因此 $J_{X^*}(f)=x^{***}$, 故 $J_{X^*}$ 是满射, 从而 $X^*$ 自反.

{\heiti 充分性:} 假设 $X^*$ 自反, 即 $J_{X^*}$ 是满射. 因为 $J_X$ 是等距嵌入且 $X$ 完备, 所以由\cref{proposition:值域闭等价完备}知$J_X(X)$ 是 $X^{**}$ 的闭子空间. 若 $X$ 不自反, 则 $J_X(X)$ 是 $X^{**}$ 的真闭子空间，进而存在$x_0\in X^{**}\setminus J_X(X)$,使得$\rho(x_0,J_X(X))>0$。 由 \cref{theorem:泛函分析--零值泛函延拓定理}, 存在非零 $x^{***}\in X^{***}$ 使得
\begin{align*}
x^{***}|_{J_X(X)}=0.
\end{align*}
因为 $X^*$ 自反, 存在 $f\in X^*$ 使得 $J_{X^*}(f)=x^{***}$. 对任意 $x\in X$, 有
\begin{align*}
0=x^{***}(J_Xx)=J_{X^*}(f)(J_Xx)=(J_Xx)(f)=f(x).
\end{align*}
故 $f=0$, 从而 $x^{***}=0$, 与 $x^{***}$ 非零矛盾! 因此 $X$ 自反.
\end{proof}

\begin{example}
在$l^1$中定义算子
\begin{align*}
T:(x_1,x_2,\cdots,x_n,\cdots)\mapsto(0,x_1,x_2,\cdots,x_n,\cdots),
\end{align*}
求证: $T\in\mathscr{L}(l^1)$并求$T^*$.
\end{example}

\begin{example}
在$l^2$中定义算子
\begin{align*}
T:(x_1,x_2,\cdots,x_n,\cdots)\mapsto\left(x_1,\frac{x_2}{2},\cdots,\frac{x_n}{n},\cdots\right),
\end{align*}
求证: $T\in\mathscr{L}(l^2)$并求$T^*$.
\end{example}

\begin{definition}[Hilbert 空间伴随算子]
设 \(H\) 是Hilbert 空间，内积记为 \((\cdot,\cdot)\),再设\(A \in \mathscr{L}(H)\)。
若存在有界线性算子 \(A^*: H \to H\)，使得对任意的 \(x, y \in H\)，都有
\begin{align*}
(Ax,y) = (x, A^*y)
\end{align*}
成立，则称 \(A^*\) 为 \(A\) 的 \textbf{Hilbert空间上的伴随算子}。特别的,若$A=A^*$,则称$A$为\textbf{Hilbert空间上的自伴随算子}.
\end{definition}

\begin{proposition}\label{proposition:Hilbert空间有界线性算子必要唯一内积伴随的共轭算子}
设 $H$ 是 Hilbert 空间, $A\in\mathscr{L}(H)$. 则存在唯一的Hilbert空间上的伴随算子 $A^*\in\mathscr{L}(H)$.进而
\begin{align*}
(Ax,y)=(x,A^*y)\quad(\forall x,y\in H).
\end{align*}
\end{proposition}
\begin{proof}

定义\(J:H\to H^*\)为
\begin{align}\label{eq:Jfhjwijefiohuie33392}
J(y)(x)=(x,y),\quad \forall x,y\in H.
\end{align}
显然$J$是共轭线性的.
由\hyperref[theorem:Riesz表示定理(Hilbert空间)]{Riesz表示定理}知,对$\forall f\in H^*$,存在唯一的$y_f\in H$,使得$\|f\|=\|y_f\|$,并且
\begin{align*}
f(x) = (x,y_f) =J(y_f)(x)\,\,(\forall x\in H)\Longrightarrow f=J(y_f).
\end{align*}
故\(J\)是双射.又由$\|f\|=\|y_f\|$可得
\begin{align*}
\|J(y_f)\| = \underset{\left\| x \right\| =1}{\mathrm{sup}}\left| J\left( y_f \right) \left( x \right) \right|=\underset{\left\| x \right\| =1}{\mathrm{sup}}\left| \left( x,y_f \right) \right|=\underset{\left\| x \right\| =1}{\mathrm{sup}}\left| f\left( x \right) \right|=\left\| f \right\| .
\end{align*}
因此\(J\)共轭线性等距同构.

由\cref{proposition:Banach空间共轭算子必存在唯一}可设\(A':H^*\to H^*\) 为\(A\) 的共轭算子,则\(A'\) 满足
\begin{align}\label{eq:Jfhjwijefiohuie33392-1}
(A'f)(x)=f(Ax),\quad \forall f\in H^*, x\in H.
\end{align}
再定义$A^*:H\to H$为
\begin{align*}
A^*=J^{-1} A' J.
\end{align*}
因为 \(J\), \(J^{-1}\) 和 \(A'\) 都是有界线性算子，所以\(A^*\) 仍是有界线性算子。
对任意 \(x, y\in H\),由\eqref{eq:Jfhjwijefiohuie33392}\eqref{eq:Jfhjwijefiohuie33392-1}式知
\begin{align*}
(Ax,y)=J(y)(Ax)=(A^{\prime}J(y))(x)=J\bigl( J^{-1}A^{\prime}J(y) \bigr) (x)=J(A^*y)(x)=(x,A^*y).
\end{align*}
这便证明了存在性。

下证唯一性.若 $B\in\mathscr{L}(H)$ 也满足
\begin{align*}
(Ax,y)=(x,By)\quad(\forall x,y\in H),
\end{align*}
则
\begin{align*}
(x,A^*y-By)=0\quad(\forall x\in H).
\end{align*}
取 $x=A^*y-By$, 得 $\|A^*y-By\|^2=0$, 故 $A^*y=By$. 由 $y$ 的任意性得 $B=A^*$.
\end{proof}

\begin{example}
设$H$是Hilbert空间，$A\in\mathscr{L}(H)$并满足
\begin{align*}
(Ax,y)=(x,Ay)\quad(\forall x,y\in H),
\end{align*}
求证:
\begin{enumerate}[(1)]
\item $A$是Hilbert空间上的自伴随算子;
\item 若$R(A)$在$H$中稠密，则方程$Ax=y$对$\forall y\in R(A)$存在唯一解.
\end{enumerate}
\end{example}
\begin{proof}
\begin{enumerate}[(1)]
\item 由\cref{proposition:Hilbert空间有界线性算子必要唯一内积伴随的共轭算子}知,存在唯一的Hilbert空间上的伴随算子$A^*\in \mathscr{L}(H)$,使得
\begin{align*}
(Ax,y)=(x,A^*y)\quad(\forall x,y\in H).
\end{align*}
又由题设知
\begin{align*}
(Ax,y)=(x,Ay)\quad(\forall x,y\in H).
\end{align*}
故$(x,A^*y-Ay)$对任意$x,y\in H$都成立.再取$x=A^*y-Ay$,则$\|A^*y-Ay\|^2=0$,
由内积的正定性可得 $A^*y=Ay$, 故 $A^*=A$.

\item 先证明 $\ker A=R(A)^\perp$. 若 $x\in\ker A$, 则对任意 $u\in H$,
\begin{align*}
(x,Au)=(Ax,u)=0,
\end{align*}
所以 $x\in R(A)^\perp$. 反过来, 若 $x\in R(A)^\perp$, 则对任意 $u\in H$,有
\begin{align*}
(Ax,u)=(x,Au)=0.
\end{align*}
取 $u=Ax$, 得 $(Ax,Ax)=0$, 于是 $Ax=\theta$, 即 $x\in\ker A$. 因此 $\ker A=R(A)^\perp$.

由于 $R(A)$ 在 $H$ 中稠密, 故由\subcref{proposition:Hilbert空间稠密线性空间正交补为0}{proposition:Hilbert空间稠密线性空间正交补为0-2}知$R(A)^\perp=\{\theta\}$, 从而 $\ker A=\{\theta\}$, 即 $A$ 是单射.
故方程$Ax=y$对$\forall y\in R(A)$存在唯一解.
\end{enumerate}
\end{proof}

\begin{example}
设$\mathscr{X},\mathscr{Y}$是$B^*$空间，$A\in\mathscr{L}(\mathscr{X},\mathscr{Y})$，又设$A^{-1}$存在且$A^{-1}\in\mathscr{L}(\mathscr{Y},\mathscr{X})$，求证:
\begin{enumerate}[(1)]
\item $(A^*)^{-1}$存在，且$(A^*)^{-1}\in\mathscr{L}(\mathscr{X}^*,\mathscr{Y}^*)$;
\item $(A^*)^{-1}=(A^{-1})^*$.
\end{enumerate}
\end{example}
\begin{proof}

因为$A^{-1}\in\mathscr{L}(\mathscr{Y},\mathscr{X})$，所以其共轭算子$(A^{-1})^*\in\mathscr{L}(\mathscr{X}^*,\mathscr{Y}^*)$。下面验证$(A^{-1})^*$是$A^*$的逆算子。

任取$f\in\mathscr{X}^*$和$x\in\mathscr{X}$。由共轭算子的定义，
\begin{align*}
\bigl(A^*((A^{-1})^*f)\bigr)(x)
=((A^{-1})^*f)(Ax)
=f(A^{-1}(Ax))
=f(x).
\end{align*}
因此$A^*((A^{-1})^*f)=f$，即
\begin{align*}
A^*\circ(A^{-1})^*=I_{\mathscr{X}^*}.
\end{align*}
任取$g\in\mathscr{Y}^*$和$y\in\mathscr{Y}$。同样地，
\begin{align*}
\bigl((A^{-1})^*(A^*g)\bigr)(y)
=(A^*g)(A^{-1}y)
=g(A(A^{-1}y))
=g(y).
\end{align*}
因此$(A^{-1})^*(A^*g)=g$，即
\begin{align*}
(A^{-1})^*\circ A^*=I_{\mathscr{Y}^*}.
\end{align*}
综上，$A^*$可逆，且$(A^*)^{-1}=(A^{-1})^*$。由于$(A^{-1})^*\in\mathscr{L}(\mathscr{X}^*,\mathscr{Y}^*)$，故$(A^*)^{-1}\in\mathscr{L}(\mathscr{X}^*,\mathscr{Y}^*)$。
\end{proof}

\begin{example}
设$\mathscr{X},\mathscr{Y},\mathscr{Z}$是$B^*$空间，而$B\in\mathscr{L}(\mathscr{X},\mathscr{Y})$以及$A\in\mathscr{L}(\mathscr{Y},\mathscr{Z})$，求证: $(AB)^*=B^*A^*$.
\end{example}

\begin{example}
设$\mathscr{X},\mathscr{Y}$是$B$空间，$T$是$\mathscr{X}$到$\mathscr{Y}$的线性算子，又设对$\forall g\in\mathscr{Y}^*$，$g(Tx)$是$\mathscr{X}$上的有界线性泛函，求证: $T$是连续的.
\end{example}

\begin{example}
设$\{x_n\}\subseteq C[a,b]$，$x\in C[a,b]$且$x_n\to x\ (n\to\infty)$，求证:
\begin{align*}
\lim\limits_{n\to\infty}x_n(t)=x(t)\quad(\forall t\in[a,b])\quad(\text{点点收敛}).
\end{align*}
\end{example}

\begin{example}
已知在$B^*$空间中$x_n\rightharpoonup x_0\ (n\to\infty)$，求证:
\begin{align*}
\varliminf\limits_{n\to\infty}\|x_n\|\geqslant\|x_0\|.
\end{align*}
\end{example}

\begin{example}
设$H$是Hilbert空间，$\{e_n\}$是$H$的正交规范基，求证: 在$H$中$x_n\rightharpoonup x_0\ (n\to\infty)$的充要条件是
\begin{enumerate}[(1)]
\item $\|x_n\|$有界;
\item $(x_n,e_k)\to(x_0,e_k)\ (n\to\infty)(k=1,2,\cdots)$.
\end{enumerate}
\end{example}

\begin{example}
设$S_n$是$L^p(\mathbb{R})(1\leqslant p<\infty)$到自身的算子:
\begin{align*}
(S_nu)(x)=
\begin{cases}
u(x),&|x|\leqslant n,\\
0,&|x|>n,
\end{cases}
\end{align*}
其中$u\in L^p(\mathbb{R})$是任意的，求证: $\{S_n\}$强收敛于恒同算子$I$，但不一致收敛到$I$.
\end{example}

\begin{example}
设$H$是Hilbert空间，在$H$中$x_n\rightharpoonup x_0\ (n\to\infty)$，而且$y_n\rightharpoonup y_0\ (n\to\infty)$，求证: $(x_n,y_n)\to(x_0,y_0)\ (n\to\infty)$.
\end{example}

\begin{example}
设$\{e_n\}$是Hilbert空间$H$中的正交规范集，求证: 在$H$中$e_n\rightharpoonup \theta\ (n\to\infty)$，但$e_n\nrightarrow \theta\ (n\to\infty)$.
\end{example}

\begin{proposition}
求证: 在自反的$B$空间中, 集合的弱列紧性与有界性是等价的.
\end{proposition}
\begin{proof}

设 $X$ 为自反的 Banach 空间, $A\subseteq X$. 

{\heiti 必要性:} 设 $A$ 弱列紧. 若 $A$ 无界, 则对每个 $n\in\mathbb{N}$, 存在 $x_n\in A$, 使得 $\|x_n\|>n$. 由弱列紧性, $\{x_n\}$ 有子列 $\{x_{n_k}\}$ 弱收敛于某 $x\in X$. 由于\cref{theorem:弱收敛的充要条件111}, 故 $\{x_{n_k}\}$ 有界; 但 $\|x_{n_k}\|>n_k\geqslant k$, 矛盾！ 故 $A$ 有界.

{\heiti 充分性:} 设 $A$ 有界. 任取 $\{x_n\}\subseteq A$. 存在 $C>0$ 使得 $\|x_n\|\leqslant C$. 令
\begin{align*}
Y=\overline{\operatorname{span}\{x_n:n\in\mathbb{N}\}}.
\end{align*}
由于 $X$ 自反, 故由\hyperref[theorem:Pettis定理]{Pettis定理}知其闭子空间 $Y$ 也自反; 又$\operatorname{span}\{x_n:n\in\mathbb{N}\}$可数,$\operatorname{span}\{x_n:n\in\mathbb{N}\}\subset Y$且$Y=\overline{\operatorname{span}\{x_n:n\in\mathbb{N}\}},$ 故 $Y$ 可分. 由\hyperref[theorem:泛函分析------Banach定理]{Banach定理}知 $Y^*$ 可分. 取 $Y^*$ 的可数稠密子集 $\{f_k:k\in\mathbb{N}\}$.

因为 $\{x_n\}$ 有界, 所以由\cref{proposition:有界线性算子的基本不等式}知,对每个 $k$, 数列 $\{f_k(x_n)\}_{n=1}^{\infty}$ 有界.利用\hyperref[proposition:对角线法则]{对角线法则},可选取子列 $\{x_{n_k}\}$, 使得对每个 $k$, $\lim\limits_{k\to\infty} f_k(x_{n_k})$都存在.

下面证明对任意 $f\in Y^*$, 数列 $\{f(x_{n_k})\}$ 收敛. 任取 $f\in Y^*$ 和 $\varepsilon>0$. 由$\{f_k:k\in\mathbb{N}\}$是$Y^*$的可数稠密子集知,存在 $k$ 使 $\|f-f_k\|<\frac{\varepsilon}{3C}$. 因 $\{f_k(x_{n_k})\}$ 收敛, 存在 $K$, 使当 $k,l\geqslant K$ 时 $|f_k(x_{n_k})-f_k(x_{n_l})|<\frac{\varepsilon}{3}$. 于是再由\cref{proposition:有界线性算子的基本不等式}可得
\begin{align*}
|f(x_{n_k})-f(x_{n_l})|
&\leqslant |f(x_{n_k})-f_k(x_{n_k})|+|f_k(x_{n_k})-f_k(x_{n_l})|+|f_k(x_{n_l})-f(x_{n_l})|\\
&\leqslant \left( \left\| x_{n_k} \right\| +\left\| x_{n_l} \right\| \right) \left\| f-f_k \right\| +\frac{\varepsilon}{3}\leqslant 2C\left\| f-f_k \right\| +\frac{\varepsilon}{3}<\varepsilon .
\end{align*}
所以 $\{f(x_{n_k})\}$ 是 Cauchy 列, 又由\cref{theorem:共轭空间的完备性}知$Y^*$完备,故收敛.

定义 $g: Y^*\to\mathbb{K}$ 为
\begin{align*}
g(f)=\lim\limits_{k\to\infty} f(x_{n_k}),\,\,f\in Y^*.
\end{align*}
显然 $g$ 线性, 且由\cref{proposition:有界线性算子的基本不等式}知$|g(f)|\leqslant C\|f\|$, 故 $g\in Y^{**}$. 由于 $Y$ 自反, 存在 $x\in Y$, 使得 $g(f)=f(x)$ 对一切 $f\in Y^*$ 成立. 于是对任意 $f\in Y^*$, 有 $f(x_{n_k})\to f(x)$. 对任意 $F\in X^*$, 令 $f=F|_Y\in Y^*$, 则 $F(x_{n_k})=f(x_{n_k})\to f(x)=F(x)$. 因此 $x_{n_k}\rightharpoonup x$ 于 $X$. 故 $A$ 弱列紧.
\end{proof}

\begin{example}
求证: $B^*$空间中的闭凸集是弱闭的，即若$M$是闭凸集，$\{x_n\}\subseteq M$，且$x_n\rightharpoonup x_0\ (n\to\infty)$，则$x_0\in M$.
\end{example}

\begin{example}
设$\mathscr{X}$是自反的$B$空间，$M$是$\mathscr{X}$中的有界闭凸集，$\forall f\in\mathscr{X}^*$，求证: $f$在$M$上达到最大值和最小值.
\end{example}

\begin{example}
设$\mathscr{X}$是自反的$B$空间，$M$是$\mathscr{X}$中的非空闭凸集，求证: $\exists x_0\in M$，使得$\|x_0\|=\inf\limits\{\|x\|\mid x\in M\}$.
\end{example}

































\end{document}